\documentclass[11pt,oneside]{article}

% =====================================================================
%  Recurrence-Free Dyadic Values of the Fabius and Rvachev Functions
%  Consolidated document.  Editorial merge (2026-09-04) of the three
%  independently written articles filed on 2026-09-04 in the subgroup
%  inverse-and-sampling/dyadic-up-extraction/; see the provenance appendix
%  for the absorbed sources and what each contributed.
%
%  House style follows the canonical corpus volumes: A4, Libertinus with an
%  lmodern fallback, article class.  Mathematical notation is normative from
%  docs/FabiusFunction_Mathematical_Notation_Catalogue through
%  fabius-notation.tex; document-local shorthands are listed below and in
%  the notation appendix.  Labels carry the prefix rf:.
% =====================================================================

\usepackage[a4paper,margin=26mm,headheight=15pt]{geometry}
\usepackage[T1]{fontenc}
\usepackage[utf8]{inputenc}
\IfFileExists{libertinus.sty}{\usepackage{libertinus}}{\usepackage{lmodern}}
\usepackage{microtype}
\usepackage{amsmath,amssymb,amsthm,mathtools,mathrsfs}
% Canonical mathematical notation for active FabiusFunction LaTeX documents.
%
% This file is intentionally a plain TeX input rather than a LaTeX package.
% Every standalone document loads it by an explicit path relative to that
% document's directory; included fragments inherit the definitions from their
% root document.  Keep this file independent of document layout and metadata.
%
% Contract:
%   * one command has one mathematical meaning throughout the active corpus;
%   * names encode semantic roles rather than a document-local choice of letter;
%   * visually similar but differently normalized objects have different print;
%   * every command below is defined normatively in the notation catalogue.
%
% The active corpus excludes docs/archive/.  Historical sources may retain
% their original notation only inside an explicitly labelled quotation.

\makeatletter
\@ifundefined{mathbb}{%
  \PackageError{fabius-notation}{The amssymb package must be loaded first}%
    {Load amsmath and amssymb before inputting fabius-notation.tex.}%
}{}
\@ifundefined{operatorname}{%
  \PackageError{fabius-notation}{The amsmath package must be loaded first}%
    {Load amsmath before inputting fabius-notation.tex.}%
}{}
\makeatother

% Version stamp used by the catalogue and by static audits.
\newcommand{\FabiusNotationVersion}{2026-08-31}

% Number systems.  NaturalNumbers is explicitly the nonnegative convention.
\newcommand{\NaturalNumbers}{\mathbb{N}_{0}}
\newcommand{\PositiveIntegers}{\mathbb{N}_{>0}}
\newcommand{\IntegerNumbers}{\mathbb{Z}}
\newcommand{\RationalNumbers}{\mathbb{Q}}
\newcommand{\RealNumbers}{\mathbb{R}}
\newcommand{\ComplexNumbers}{\mathbb{C}}
\newcommand{\ExtendedRealNumbers}{\overline{\mathbb{R}}}

% The three principal Fabius--Rvachev functions.  Subscripts are deliberate:
% a bare F and a calligraphic F have several unrelated uses in the corpus.
\newcommand{\FabiusBounded}{F_{\mathrm{bd}}}
\newcommand{\FabiusGlobal}{F_{\mathrm{gl}}}
\newcommand{\FabiusQuantile}{Q_{F}}
\newcommand{\FabiusClampedQuantile}{Q_{F,\mathrm{cl}}}
\newcommand{\RvachevUp}{\operatorname{up}}

% Constants, elementary operators, and delimiters.
\newcommand{\EulerE}{\mathrm{e}}
\newcommand{\ImaginaryUnit}{\mathrm{i}}
\newcommand{\Differential}{\mathop{}\!\mathrm{d}}
\newcommand{\AbsoluteValue}[1]{\left\lvert #1\right\rvert}
\newcommand{\Norm}[1]{\left\lVert #1\right\rVert}
\newcommand{\Floor}[1]{\left\lfloor #1\right\rfloor}
\newcommand{\Ceiling}[1]{\left\lceil #1\right\rceil}
\newcommand{\FractionalPart}[1]{\left\{#1\right\}}
\newcommand{\PositivePart}[1]{\left(#1\right)_{+}}
\newcommand{\IndicatorOf}[1]{\mathbf{1}_{#1}}
\newcommand{\SetOf}[1]{\left\{#1\right\}}
\newcommand{\InnerProduct}[1]{\left\langle #1\right\rangle}
\newcommand{\CardinalityOf}[1]{\left\lvert #1\right\rvert}
\newcommand{\DefinitionEquals}{\mathrel{\coloneqq}}
\newcommand{\Sign}{\operatorname{sgn}}
\newcommand{\IdentityOperator}{\operatorname{id}}
\newcommand{\SpanOperator}{\operatorname{span}}
\newcommand{\TraceOperator}{\operatorname{tr}}
\newcommand{\RankOperator}{\operatorname{rank}}
\newcommand{\DiagonalOperator}{\operatorname{diag}}
\newcommand{\ResidueOperator}{\operatorname*{Res}}
\newcommand{\LeastCommonMultiple}{\operatorname{lcm}}
\newcommand{\OrderOperator}{\operatorname{ord}}
\newcommand{\PrincipalLogarithm}{\operatorname{Log}}
\newcommand{\FallingFactorial}[2]{\left(#1\right)^{\underline{#2}}}
\newcommand{\RisingFactorial}[2]{\left(#1\right)^{\overline{#2}}}

% Support, topology, convolution, and metric notation.
\newcommand{\SupportOperator}{\operatorname{supp}}
\newcommand{\SupportOf}[1]{\SupportOperator\!\left(#1\right)}
\newcommand{\TopologicalSupportOf}[1]{\operatorname{tsupp}\!\left(#1\right)}
\newcommand{\Convolution}{\mathbin{*}}
\newcommand{\MetricDistance}{\operatorname{dist}}
\newcommand{\TotalVariationDistance}{d_{\mathrm{TV}}}

% Probability.  Equality and convergence in law are not metric distance.
\newcommand{\Probability}{\mathbb{P}}
\newcommand{\Expectation}{\mathbb{E}}
\newcommand{\Variance}{\operatorname{Var}}
\newcommand{\Covariance}{\operatorname{Cov}}
\newcommand{\LawOperator}{\operatorname{Law}}
\newcommand{\LawOf}[1]{\LawOperator\!\left(#1\right)}
\newcommand{\EqualInLaw}{\mathrel{\overset{\mathrm{law}}{=}}}
\newcommand{\ConvergesInLaw}{\mathrel{\xrightarrow{\mathrm{law}}}}

% Fourier and integral transforms.  The printed labels expose normalization.
% FourierTwoPi uses exp(-2*pi*i*x*xi); FourierAngular uses exp(-i*x*omega).
\newcommand{\FourierTwoPi}[1]{\widehat{#1}_{\,2\pi}}
\newcommand{\FourierAngular}[1]{\widehat{#1}_{\,\mathrm{ang}}}
\newcommand{\LaplaceTransformSymbol}{\mathcal{L}}
\newcommand{\MellinTransformSymbol}{\mathcal{M}}
\newcommand{\LaplaceTransformOf}[1]{\LaplaceTransformSymbol\!\left[#1\right]}
\newcommand{\MellinTransformOf}[1]{\MellinTransformSymbol\!\left[#1\right]}
\newcommand{\MomentGeneratingFunction}[1]{M_{#1}}
\newcommand{\CharacteristicFunction}[1]{\varphi_{#1}}

% The two standard sinc functions must never share one symbol.
\newcommand{\SincRad}{\operatorname{sinc}_{\mathrm{rad}}}
\newcommand{\SincPi}{\operatorname{sinc}_{\pi}}

% Binary arithmetic and Thue--Morse notation.
\newcommand{\BinaryDigitSum}{\operatorname{wt}_{2}}
\newcommand{\ThueMorseSignSymbol}{\varepsilon}
\newcommand{\ThueMorseSign}[1]{\varepsilon_{#1}}
\newcommand{\PadicValuation}[1]{\nu_{#1}}
\newcommand{\TwoAdicValuation}{\nu_{2}}

% q-series notation.  Every base is an explicit argument.
\newcommand{\QPochhammer}[3]{\left(#1;#2\right)_{#3}}
\newcommand{\GaussianBinomial}[3]{%
  \genfrac{[}{]}{0pt}{}{#1}{#2}_{#3}%
}
\newcommand{\QInteger}[2]{\left[#1\right]_{#2}}

% Named special functions.
\newcommand{\Polylogarithm}{\operatorname{Li}}
\newcommand{\LambertW}{W}
\newcommand{\GammaFunction}{\Gamma}
\newcommand{\RiemannZeta}{\zeta}

% Asymptotics.  BigO and LittleO are operator tokens for existing formula
% grammar; prefer the At forms so that the limiting regime is printed.
\newcommand{\BigO}{\mathcal{O}}
\newcommand{\LittleO}{o}
\newcommand{\BigOAt}[2]{\mathcal{O}_{#2}\!\left(#1\right)}
\newcommand{\LittleOAt}[2]{o_{#2}\!\left(#1\right)}

\usepackage{booktabs,longtable,tabularx,array}
\usepackage{xcolor}
\usepackage{enumitem}
\usepackage{fancyhdr}
\usepackage{xurl}
\usepackage{listings}
\usepackage{hyperref}
\usepackage[nameinlink,noabbrev,capitalise]{cleveref}

\definecolor{linkblue}{RGB}{24,72,124}
\definecolor{deepgreen}{RGB}{23,102,69}
\definecolor{shadegray}{RGB}{246,247,249}
\definecolor{rulegray}{RGB}{195,201,208}

\hypersetup{
  hypertexnames=false,
  colorlinks=true,
  linkcolor=linkblue,
  citecolor=deepgreen,
  urlcolor=linkblue,
  pdftitle={Recurrence-Free Dyadic Values of the Fabius and Rvachev Functions},
  pdfauthor={Vladimir Reshetnikov; ProveIt formal-development synthesis},
  pdfsubject={Consolidated document: finite, recurrence-free evaluation of the Fabius distribution function and Rvachev's up-function at dyadic rationals},
  pdfkeywords={Fabius function, Rvachev up-function, dyadic rational, Thue-Morse sequence, Bernoulli numbers, ordered compositions, finite convolution, Gaussian binomial coefficient, exact extrapolation}
}

\pagestyle{fancy}
\fancyhf{}
\fancyhead[L]{\small Recurrence-free dyadic values}
\fancyhead[R]{\small\nouppercase{\leftmark}}
\fancyfoot[C]{\thepage}
\renewcommand{\headrulewidth}{0.4pt}

\setcounter{tocdepth}{2}
\setcounter{secnumdepth}{3}
\setlist[itemize]{leftmargin=1.45em,itemsep=0.25ex,topsep=0.55ex}
\setlist[enumerate]{leftmargin=1.85em,itemsep=0.35ex,topsep=0.55ex}
\allowdisplaybreaks[2]
\setlength{\emergencystretch}{2.5em}

\theoremstyle{plain}
\newtheorem{theorem}{Theorem}[section]
\newtheorem{proposition}[theorem]{Proposition}
\newtheorem{lemma}[theorem]{Lemma}
\newtheorem{corollary}[theorem]{Corollary}
\theoremstyle{definition}
\newtheorem{definition}[theorem]{Definition}
\newtheorem{example}[theorem]{Example}
\theoremstyle{remark}
\newtheorem{remark}[theorem]{Remark}

\lstdefinestyle{pseudo}{%
  basicstyle=\ttfamily\small,columns=fullflexible,frame=single,
  rulecolor=\color{rulegray},backgroundcolor=\color{shadegray},
  showstringspaces=false,keepspaces=true,breaklines=true,
  xleftmargin=0.4em,xrightmargin=0.4em,aboveskip=0.8em,belowskip=0.8em}
\lstset{style=pseudo}

% ---------------------------------------------------------------------
%  Document-local shorthands.  Each expands to catalogue notation or to an
%  object defined in the setting section; recorded in the notation appendix.
% ---------------------------------------------------------------------
\newcommand{\Fb}{\FabiusBounded}            % bounded distribution function
\newcommand{\Fg}{\FabiusGlobal}             % signed global extension
\newcommand{\up}{\RvachevUp}                % the even bump
\newcommand{\tm}[1]{\ThueMorseSign{#1}}     % Thue--Morse sign
\newcommand{\E}{\Expectation}
\newcommand{\Prob}{\Probability}
\newcommand{\dd}{\Differential}
\newcommand{\coeff}[2]{\CoefficientExtraction{#1}{#2}}
\newcommand{\qb}[3]{\GaussianBinomial{#1}{#2}{#3}}
\newcommand{\qp}[3]{\QPochhammer{#1}{#2}{#3}}
\newcommand{\Comp}[1]{\mathcal{C}_{#1}}     % ordered compositions of #1
\newcommand{\Prof}[1]{\mathcal{P}_{#1}}     % multiplicity profiles of #1
\newcommand{\Spow}[2]{\mathcal{S}_{#1}(#2)} % integer Thue--Morse power sum
\newcommand{\Tr}[1]{T_{#1}}                 % triangular number
\newcommand{\Dop}{\mathrm{D}}               % differentiation of polynomials
\newcommand{\amvol}[1]{a^{\mathrm{vol}}_{#1}} % the canonical volume's reciprocal coefficients

\begin{document}
\title{\textbf{Recurrence-Free Dyadic Values}\\[0.35em]
  \Large of the Fabius and Rvachev Functions\\[0.6em]
  \large Finite moment elimination, binary compression, integer determinants,\\
  and exact quarter-base spline extraction}
\author{Vladimir Reshetnikov\\
  \small ProveIt formal-development synthesis}
\date{4 September 2026}

\maketitle

\begin{abstract}
Exact rationality of the Fabius distribution function $\Fb$ and of Rvachev's
up-function at dyadic arguments is classical, but the formulas that display
it usually leave an auxiliary sequence --- the even moments, or the values at
reciprocal powers of two --- defined only by a triangular recurrence, or
reach the value as the limit of finite convolutions.  This document removes
both dependencies and organizes every finite route to the value in one
place.

The principal formula expresses $\Fb(m/2^{n})$ through
$\Floor{n/2}+1$ explicit Thue--Morse power sums with rational
interpolation weights at the quarter base $q=\tfrac14$.  Behind it is an
exact, not asymptotic, statement: at a fixed dyadic point of depth $n$ the
centred finite-uniform distribution functions are polynomials of degree at
most $\Floor{n/2}$ in $4^{-N}$, and a shrinking-cell version identifies the
whole polynomial piece of the finite spline with a finite deconvolution of
the dyadic Taylor polynomial.  The degree is exactly $\Floor{n/2}$ at every
reduced interior dyadic, so the sample count is optimal for any linear rule
whose weights work on the whole grid.

A second family compresses the classical Thue--Morse moment formula to one
outer summand per nonzero binary digit of the numerator --- in two different
ways, a block identity and a telescope --- and supplies the coefficients by
positive ordered-composition sums, by Bernoulli multiplicity sums with the
Bernoulli numbers themselves given as finite sums, and by interpolation of
finitely many uniform convolutions.  An integer bordered determinant packages
the value as a single ratio of integers and yields an explicit common
denominator, which in turn gives a certified rounding formula.  The same
finite algebra evaluates the bump, the signed global extension, every dyadic
derivative, all iterated primitives, and the analogous normalized
reciprocal-integer-base laws.  Every result carries a complete proof; every
printed rational is checked by the retained exact-arithmetic verifier; the
relation to the canonical extraction volume of this subgroup is stated
formula by formula.  No Lean formalization of the new material is claimed.
\end{abstract}

\vfill
\noindent\textbf{Scope.}  The bounded distribution function, the compactly
supported bump, and the signed global extension are kept distinct
throughout.  The results are mathematical derivations with complete proofs;
the formalization register at the end records which pieces have Lean
counterparts in the repository and which do not.

\clearpage
\tableofcontents
\clearpage

% =====================================================================
\section{The question and the map of formulas}
\label{rf:sec:intro}

\subsection{What counts as a closed form here}
\label{rf:sec:intro-closed}

We seek the exact values of the Fabius distribution function and of
Rvachev's bump at $m/2^{n}$, with integer $m$ and nonnegative integer $n$.
A formula is called \emph{recurrence-free} when its right-hand side
contains only explicitly bounded finite sums and products of specified
numbers, together with rational arithmetic, integer powers, factorials,
binary digits, and finite $q$-Pochhammer symbols.  A coefficient-extraction
symbol may abbreviate a finite sum that is displayed separately, and a
determinant is admissible because its entries are given explicitly and its
Leibniz expansion is a finite sum.  What is excluded is a definition that
leaves a desired auxiliary quantity specified only in terms of earlier
unknown quantities: a formula that refers to ``the moments $c_{k}$,
computed recursively'' is not a closed form in this sense, and neither is a
limit of explicit rational approximations.  Recursion may of course be an
efficient \emph{implementation}; the point is that no recurrence is needed
to \emph{define} the numbers on the right-hand side of an evaluation
identity.

Three functions must be distinguished.  Throughout, $\Fb$ is the bounded
Fabius distribution function, $\up$ is the compactly supported bump, and
$\Fg$ is the signed global extension.  On their principal domains,
\begin{equation}
 \up(x)=\Fb(1-\AbsoluteValue{x}),\qquad
 \Fb(x)=\up(x-1)\quad(0\le x\le1),\qquad
 \Fg=\Fb\ \text{on }[0,1].
 \label{rf:eq:intro-fold}
\end{equation}
The precise conventions and the probability model are fixed in
\cref{rf:sec:setup}.

\subsection{The principal formula}
\label{rf:sec:intro-principal}

Write
\begin{equation}
 \Tr{r}\DefinitionEquals\frac{r(r+1)}{2},\qquad
 d\DefinitionEquals\Floor{n/2},\qquad
 \tm{h}\DefinitionEquals(-1)^{\BinaryDigitSum(h)},
 \label{rf:eq:intro-notation}
\end{equation}
where $\BinaryDigitSum(h)$ is the number of ones in the binary expansion of
$h$.  The Thue--Morse sign is a digit formula, not a recursively defined
sequence: for $0\le h<2^{L}$ its digits are the explicit integers
$\Floor{h/2^{\ell}}-2\Floor{h/2^{\ell+1}}$, $0\le\ell<L$, and
\begin{equation}
 \tm{h}=\prod_{\ell=0}^{L-1}
 \Bigl(1-2\bigl(\Floor{h/2^{\ell}}-2\Floor{h/2^{\ell+1}}\bigr)\Bigr).
 \label{rf:eq:intro-tm-digits}
\end{equation}
Define the \emph{integer Thue--Morse power sum}
\begin{equation}
 \Spow{N}{A}\DefinitionEquals\sum_{h=0}^{A-1}\tm{h}\,(2A-2h-1)^{N}
 \qquad(N\ge1,\ A\ge0).
 \label{rf:eq:intro-power-sum}
\end{equation}
Empty sums are zero, empty products are one, and $0^{0}=1$ in the finite
polynomial formulas.

\begin{theorem}[Quarter-base closed formula]
\label{rf:thm:main-quarter}
Let $n\ge1$, $0\le m\le2^{n}$, and $d=\Floor{n/2}$.  Put
\begin{equation}
 w_{d,j}\DefinitionEquals
 \frac{(-1)^{d-j}\,4^{\Tr{j}}}
      {\displaystyle\prod_{r=1}^{j}(4^{r}-1)\prod_{r=1}^{d-j}(4^{r}-1)}
 \qquad(0\le j\le d).
 \label{rf:eq:main-weights-integer}
\end{equation}
Then
\begin{equation}
 \boxed{\;
 \Fb\!\left(\frac{m}{2^{n}}\right)
 =\sum_{j=0}^{d}w_{d,j}\,
   \frac{\Spow{n+j}{2^{j}m}}{2^{\Tr{n+j}}\,(n+j)!}\,.\;}
 \label{rf:eq:main-quarter}
\end{equation}
More generally, for every integer $R\ge0$,
\begin{equation}
 \Fb\!\left(\frac{m}{2^{n}}\right)
 =\sum_{j=0}^{d}w_{d,j}\,
   \frac{\Spow{n+R+j}{2^{R+j}m}}{2^{\Tr{n+R+j}}\,(n+R+j)!}\,.
 \label{rf:eq:main-quarter-shift}
\end{equation}
\end{theorem}

The proof is completed in \cref{rf:sec:extraction}
(\cref{rf:thm:extraction-consecutive}).  It does \emph{not} assume
rationality of the target value, so \cref{rf:eq:main-quarter} is itself a
proof of dyadic rationality using only rational finite expressions.  The
formula contains no moments, Bernoulli numbers, derivative values,
recursively defined coefficients, or limits.  The first rows of weights are
\begin{equation}
\begin{array}{c|l}
 d & (w_{d,0},\ldots,w_{d,d})\\\hline
 0 & (1)\\
 1 & (-1,\,4)/3\\
 2 & (1,\,-20,\,64)/45\\
 3 & (-1,\,84,\,-1344,\,4096)/2835
\end{array}
\label{rf:eq:main-weight-rows}
\end{equation}
Each row sums to one.  These are exact interpolation rows, not fitted
numerical coefficients.

\subsection{The families of formulas}
\label{rf:sec:intro-map}

There are several genuinely different finite routes to the same rational
number, and they are useful in different situations.
\Cref{rf:tab:intro-map} lists them with their characteristic ingredient.

\begin{table}[htbp]
\centering\small
\begin{tabularx}{\textwidth}{@{}p{0.24\textwidth}Xp{0.19\textwidth}@{}}
\toprule
Family & Finite ingredients and characteristic feature & Main result\\
\midrule
Bernoulli partitions & multiplicity vectors and finite Bernoulli power sums; visible cumulant structure & \cref{rf:thm:partition-value}\\
Positive compositions & ordered positive parts and Mersenne-type factors $4^{S}-1$; every coefficient term is positive & \cref{rf:thm:composition-value}\\
Binary blocks & one outer summand per nonzero binary digit; centred coefficients; valid for the signed extension at every numerator & \cref{rf:thm:binary-block}\\
Binary telescope & the same digits, an alternating telescope of raw-moment polynomials & \cref{rf:thm:binary-telescope}\\
Finite uniform prefixes & moments of finitely many uniforms, interpolated at the quarter base & \cref{rf:thm:moment-extraction}, \cref{rf:thm:finite-cube}\\
Integer determinant & a bordered matrix of order $\Floor{n/2}+1$; integer certificate and a common denominator & \cref{rf:thm:determinant}\\
Certified rounding & one enormous prefix, rounded to the common denominator & \cref{rf:cor:rounding}\\
Quarter-base extraction & $\Floor{n/2}+1$ centred spline levels; no auxiliary coefficients at all & \cref{rf:thm:main-quarter}\\
Direct bump extraction & the same weights applied to finite densities & \cref{rf:thm:up-quarter}\\
\bottomrule
\end{tabularx}
\caption{The families of recurrence-free formulas.  All rows are exact.}
\label{rf:tab:intro-map}
\end{table}

For a single value with a sparse binary expansion, one of the two binary
formulas usually avoids the largest sums.  For many values on one grid, the
master formula with centred coefficients permits the same explicitly
computed coefficients to be reused.  The composition expression is
transparent when positive terms and elementary denominators are wanted; the
partition expression connects to cumulants and suits symbolic manipulation;
the determinant is a compact integer certificate with a visible denominator.
The quarter-base formula is the most self-contained answer when \emph{no
auxiliary coefficient sequence at all} is acceptable: it contains only a
short geometric interpolation row and finite power sums.
\Cref{rf:sec:conclusion} compares them further.

\subsection{Relation to the canonical volume and to the literature}
\label{rf:sec:intro-relation}

This document belongs to the subgroup whose canonical volume
\cite{proveit-due} proves that, at a dyadic point $x$ of depth $s$, the
finite sinc-product spline $\up_{n}$ obeys
$\up_{n}(x)=\up(x)+\sum_{r=1}^{\Floor{s/2}}C_{r}(x)4^{-rn}$ exactly for all
$n\ge s$, and derives from it the quarter-base Gaussian-binomial extraction
row.  The present document proves the same polynomiality in the
distribution-function normalization
(\cref{rf:thm:local-cell,rf:cor:scale-polynomial}) and reads the row as an
\emph{evaluation formula}, \cref{rf:eq:main-quarter}, in which the spline
values themselves are replaced by explicit integer sums.  Everything else
here --- the moment-free coefficients, the binary compressions, the
determinant, the denominators, the shifted half-base identity, the
primitives, and the integer-base extension --- is complementary to the
volume.  \Cref{rf:sec:extraction-volume} gives the dictionary between the two
normalizations, and \cref{rf:sec:register} the formalization register.

The classical background is Arias de Reyna's arithmetic and analytic papers
\cite{arias1982,arias2018}, which contain the dyadic moment formula and the
sharp denominator theory; Haugland \cite{haugland} gives an exact dyadic
evaluator; Reshetnikov's 2019 question \cite{reshetnikov2019} posed the
half-base $q$-binomial expression recovered in \cref{rf:sec:halfbase}; and
Pinelis \cite{pinelis2020} gave a finite-convolution representation followed
by a limit, which is exactly the limiting step that the quarter-base row
removes.  The repository's primary exposition \cite{proveit-primary} already
contains the Thue--Morse moment formula, the reciprocal-power path and
composition descriptions, the Bell and cumulant machinery, and the half-base
layer; the Exponents volume \cite{proveit-exponents} contains the asymptotic
quarter-base Richardson combination.  None of those is claimed as new.  What
this document adds is the completion of every coefficient into a finite
expression, the organization of the distinct evaluation routes with full
proofs, and the exact centred cutoff identity with its precise degree.

\subsection{Conventions}
\label{rf:sec:intro-conventions}

Throughout, $n,m,N$ are nonnegative integers unless another range is
stated, and $\NaturalNumbers=\SetOf{0,1,2,\ldots}$.  For $k\ge0$, $\Comp{k}$
is the finite set of ordered compositions of $k$ into positive parts, the
empty composition being the only element of $\Comp{0}$; there are $2^{k-1}$
compositions of $k\ge1$, obtained by choosing cuts among $k-1$ gaps.  For
$\mathbf r=(r_{1},\ldots,r_{\ell})\in\Comp{k}$ the \emph{suffix sums} are
\begin{equation}
 R_{i}\DefinitionEquals r_{i}+r_{i+1}+\cdots+r_{\ell}\qquad(1\le i\le\ell),
 \label{rf:eq:intro-suffix}
\end{equation}
so that $R_{1}=k$ and the $R_{i}$ are distinct integers in $[1,k]$.  The
\emph{multiplicity-profile set} is
\begin{equation}
 \Prof{k}\DefinitionEquals
 \SetOf{(v_{1},\ldots,v_{k})\in\NaturalNumbers^{k}:
   v_{1}+2v_{2}+\cdots+kv_{k}=k},
 \label{rf:eq:intro-profiles}
\end{equation}
which for $k=0$ consists of the empty tuple.  For $0<q<1$,
\begin{equation}
 \qp{q}{q}{r}=\prod_{j=1}^{r}(1-q^{j}),\qquad
 \qb{d}{j}{q}=\frac{\qp{q}{q}{d}}{\qp{q}{q}{j}\qp{q}{q}{d-j}},
 \label{rf:eq:intro-q}
\end{equation}
finite products; at $q=\tfrac14$ every such value is rational.  No infinite
$q$-product is needed to evaluate any dyadic value in this document.
Bernoulli numbers use the convention $B_{1}=-\tfrac12$.  A table of the
symbols is in \cref{rf:sec:notation}.

% =====================================================================
\section{Normalization and the probability model}
\label{rf:sec:setup}

\subsection{The bounded distribution function and its centred density}
\label{rf:sec:setup-model}

Let $U_{1},U_{2},\ldots$ be independent random variables, each uniform on
$[0,1]$, and define
\begin{equation}
 X\DefinitionEquals\sum_{j=1}^{\infty}2^{-j}U_{j},\qquad
 \Fb(x)\DefinitionEquals\Prob(X\le x).
 \label{rf:eq:setup-X}
\end{equation}
The series converges absolutely for every outcome, since $0\le U_{j}\le1$
and $\sum_{j}2^{-j}=1$, and takes values in $[0,1]$.  Let $V_{j}=2U_{j}-1$,
uniform on $[-1,1]$, and
\begin{equation}
 Z\DefinitionEquals2X-1=\sum_{j=1}^{\infty}2^{-j}V_{j}.
 \label{rf:eq:setup-Z}
\end{equation}
The density of $Z$ is denoted $\up$; this is the standard probability
construction of the two functions \cite{arias1982,arias2018}, and the
catalogue entry \texttt{rvachev-up} normalizes $\up$ in exactly this way.

\begin{lemma}[Smoothness and normalization]
\label{rf:lem:setup-normalization}
The density $\up$ is even, nonnegative, $C^{\infty}$, and supported on
$[-1,1]$.  Moreover
\begin{align}
 \Fb'(x)&=2\up(2x-1),\label{rf:eq:setup-F-density}\\
 \up(x)&=\Fb(1-\AbsoluteValue{x}),\label{rf:eq:setup-fold}\\
 \Fb(1-x)&=1-\Fb(x),\label{rf:eq:setup-reflection}\\
 \up'(x)&=2\up(2x+1)-2\up(2x-1),\label{rf:eq:setup-up-dilation}\\
 \up(0)&=1,\qquad0\le\up\le1.\label{rf:eq:setup-up-bound}
\end{align}
The function $\Fb$ is $C^{\infty}$ on $\RealNumbers$, equals $0$ for
$x\le0$ and $1$ for $x\ge1$, is $2$-Lipschitz, and all its positive-order
derivatives vanish at $0$ and at $1$.  All derivatives of $\up$ vanish at
$\pm1$.
\end{lemma}

\begin{proof}
The law of $Z$ is symmetric because every $V_{j}$ has a symmetric law, and
it has no atoms: conditioning on all coordinates but $U_{1}$ leaves a
translated uniform distribution.  Its characteristic function is
\[
 \E\,\EulerE^{\ImaginaryUnit tZ}
 =\prod_{j\ge1}\frac{\sin(t/2^{j})}{t/2^{j}}.
\]
For any fixed positive integer $L$, bound the first $L$ factors by
$\min(1,2^{j}/\AbsoluteValue{t})$ and all remaining factors by one: this
gives decay $\BigOAt{(1+\AbsoluteValue{t})^{-L}}{L}$.  Hence the
characteristic function multiplied by any fixed power of $t$ is integrable,
and Fourier inversion with differentiation under the integral gives a
$C^{\infty}$ density.  Nonnegativity and the support $[-1,1]$ follow from
the law of the bounded sum; the affine change $Z=2X-1$ gives
\cref{rf:eq:setup-F-density}.

Replacing every $U_{j}$ by $1-U_{j}$ shows $X\EqualInLaw1-X$, which is
\cref{rf:eq:setup-reflection}.  For \cref{rf:eq:setup-fold} use the
independent-copy decomposition $Z\EqualInLaw(V+Z')/2$, with $V$ uniform on
$[-1,1]$ and $Z'$ an independent copy of $Z$.  The density of $V/2$ is one on
$[-\tfrac12,\tfrac12]$, so
\begin{equation}
 \up(x)=\Prob(2x-1\le Z'\le2x+1)=\Fb(x+1)-\Fb(x).
 \label{rf:eq:setup-up-head}
\end{equation}
For $x\ge0$ the upper endpoint is at least one, so the probability equals
$\Prob(X'\ge x)=1-\Fb(x)=\Fb(1-x)$; for $x\le0$ use evenness.  The same
display applies outside the support through the constant tails of $\Fb$.
Differentiating \cref{rf:eq:setup-up-head} and using
\cref{rf:eq:setup-F-density} gives \cref{rf:eq:setup-up-dilation}.  Formula
\cref{rf:eq:setup-up-head} gives $\up(0)=\Fb(1)-\Fb(0)=1$ and $\up\le1$, and
\cref{rf:eq:setup-F-density} then bounds $\Fb'$ by two.  Smoothness together
with the constant exterior values forces every positive-order derivative of
$\Fb$ to vanish at $0$ and $1$, and the compact support of the smooth $\up$
forces all its derivatives to vanish at $\pm1$.
\end{proof}

In particular $\Fb'(x)=2\Fb(2x)$ on $[0,\tfrac12]$.  This differential
equation must not be extended to all real $x$ with the \emph{bounded} $\Fb$
on the right-hand side; the signed extension of the next subsection is the
function that satisfies it globally.

\subsection{The signed extension and the dyadic jet}
\label{rf:sec:setup-signed}

Define the locally finite sum
\begin{equation}
 \Fg(x)\DefinitionEquals\sum_{b=0}^{\infty}\tm{b}\,\up(x-2b-1).
 \label{rf:eq:setup-global}
\end{equation}
On $[2b,2b+2]$ only the $b$th term can be nonzero, so
\begin{equation}
 \Fg(x)=\tm{b}\,\up(x-2b-1)\quad(2b\le x\le2b+2),\qquad
 \Fg=\Fb\ \text{on }[0,1],\qquad \Fg=0\ \text{on }(-\infty,0].
 \label{rf:eq:setup-global-block}
\end{equation}
At the shared endpoints $x=2b$ both adjacent formulas give zero, so the
choice of block is immaterial.

\begin{lemma}[Dilation and dyadic termination]
\label{rf:lem:setup-dilation}
The function $\Fg$ is $C^{\infty}$ on $\RealNumbers$ and satisfies
\begin{equation}
 \Fg'(x)=2\Fg(2x),\qquad
 \Fg^{(r)}(x)=2^{\Tr{r}}\Fg(2^{r}x)\quad(r\ge0).
 \label{rf:eq:setup-global-derivative}
\end{equation}
At nonnegative integers and half-integers,
\begin{equation}
 \Fg(2b)=0,\qquad\Fg(2b+1)=\tm{b},\qquad
 \Fg\!\left(\frac{m}{2}\right)=\frac12\,\tm{\Floor{m/4}}\quad(m\text{ odd}).
 \label{rf:eq:setup-integer-values}
\end{equation}
At a dyadic $x=m/2^{n}\in[0,1]$,
\begin{equation}
 \Fb^{(r)}(x)=0\qquad(r\ge n+1),
 \label{rf:eq:setup-jet-termination}
\end{equation}
and likewise $\up^{(r)}(x)=0$ for $r\ge n+1$ when $\AbsoluteValue{x}\le1$ is
dyadic of depth $n$.
\end{lemma}

\begin{proof}
Flatness of $\up$ at $\pm1$ makes the locally finite sum smooth, including
at the joining points.  Insert \cref{rf:eq:setup-up-dilation} into
\cref{rf:eq:setup-global}:
\[
 \Fg'(x)=2\sum_{b\ge0}\tm{b}\bigl\{\up(2x-4b-1)-\up(2x-4b-3)\bigr\}
 =2\sum_{c\ge0}\tm{c}\,\up(2x-2c-1)=2\Fg(2x),
\]
because the digit definition gives $\tm{2b}=\tm{b}$ and
$\tm{2b+1}=-\tm{b}$, so the two terms are the even and odd summands of the
single sum.  Repeated differentiation and the identity
$\Tr{r}+r=\Tr{r+1}$ prove \cref{rf:eq:setup-global-derivative}.  The
integer values follow from the support of $\up$ and $\up(0)=1$, $\up(\pm1)=0$;
at an odd half-integer $m/2$ the argument of the relevant translate is
$\pm\tfrac12$, where $\up=\Fb(\tfrac12)=\tfrac12$ by
\cref{rf:eq:setup-fold,rf:eq:setup-reflection}, and the block index is
$\Floor{m/4}$.

If $r\ge n+1$, the argument $2^{r}m/2^{n}$ is an even integer, so
\cref{rf:eq:setup-global-derivative} gives zero.  Derivatives of $\Fb$ and
$\Fg$ agree at interior points of $[0,1]$ and, by smoothness and one-sided
agreement, at the endpoints.  On $(-1,1)$ one has $\up(x)=\Fg(x+1)$, and the
same reasoning applies to $2^{r}(x+1)$; at $x=\pm1$ all derivatives of $\up$
vanish.
\end{proof}

\begin{remark}[A terminating jet is not a local polynomial identity]
\label{rf:rmk:setup-jet}
The vanishing \cref{rf:eq:setup-jet-termination} does not assert that $\Fb$
equals its Taylor polynomial near a dyadic point --- the Fabius function is
nowhere analytic.  The exact spline statements proved in
\cref{rf:sec:splines} concern a \emph{finite convolution} and a cell whose
width shrinks with its depth, together with an exact smoothing identity; at
no point is $\Fb$ replaced by a convergent Taylor series.  This distinction
is essential to every proof below.
\end{remark}

\subsection{Moment generating functions and independent tails}
\label{rf:sec:setup-mgf}

Put
\begin{equation}
 H(t)\DefinitionEquals\frac{\sinh t}{t},\quad H(0)=1,\qquad
 M(t)\DefinitionEquals\E\,\EulerE^{tZ}=\prod_{j=1}^{\infty}H\!\left(\frac{t}{2^{j}}\right).
 \label{rf:eq:setup-M}
\end{equation}
The product converges locally uniformly, since $H(t/2^{j})=1+\BigOAt{4^{-j}}{K}$
on each compact set $K$; it defines an entire even function, which is also
the moment generating function of the bounded variable $Z$ directly.  Its
even Taylor coefficients are the normalized moments
\begin{equation}
 a_{k}\DefinitionEquals\coeff{t}{2k}M(t)=\frac{\E Z^{2k}}{(2k)!},\qquad
 c_{k}\DefinitionEquals\E Z^{2k}=(2k)!\,a_{k},
 \label{rf:eq:setup-a-c}
\end{equation}
and the odd moments vanish.  Define the finite prefixes
\begin{equation}
 Z_{N}\DefinitionEquals\sum_{j=1}^{N}2^{-j}V_{j},\qquad
 M_{N}(t)\DefinitionEquals\E\,\EulerE^{tZ_{N}}=\prod_{j=1}^{N}H\!\left(\frac{t}{2^{j}}\right),
 \label{rf:eq:setup-prefix}
\end{equation}
with $Z_{0}=0$, $M_{0}=1$, and the uncentred prefixes
$X_{N}\DefinitionEquals\sum_{j=1}^{N}2^{-j}U_{j}$.  The exact
factorizations are
\begin{equation}
 M(t)=M_{N}(t)\,M(2^{-N}t),\qquad M(2t)=H(t)\,M(t),
 \label{rf:eq:setup-tail-M}
\end{equation}
the second being the first with $N=1$ after the substitution $t\mapsto2t$,
or equivalently the statement $Z\EqualInLaw(V+Z')/2$ on the transform side.
Quotients by $M$ will be used only near zero or as formal power series with
constant term one; no assertion that $1/M$ is entire is needed.  For the
raw variable,
\begin{equation}
 M_{X}(z)\DefinitionEquals\E\,\EulerE^{zX}=\EulerE^{z/2}M(z/2),\qquad
 A_{m}\DefinitionEquals\frac{\E X^{m}}{m!}=\coeff{z}{m}M_{X}(z).
 \label{rf:eq:setup-raw}
\end{equation}
Comparing coefficients in $M_{X}(z)=\EulerE^{z/2}M(z/2)$ gives the dictionary
\begin{equation}
 A_{m}=2^{-m}\sum_{k=0}^{\Floor{m/2}}\frac{a_{k}}{(m-2k)!},
 \label{rf:eq:setup-raw-centred}
\end{equation}
so a finite formula for the $a_{k}$ is also one for the $A_{m}$, and
conversely.

% =====================================================================
\section{From a finite box to the dyadic master identity}
\label{rf:sec:master}

\subsection{A finite-convolution formula with proof}
\label{rf:sec:master-box}

\begin{lemma}[Box--simplex inclusion--exclusion]
\label{rf:lem:master-box}
Let $w_{1},\ldots,w_{N}>0$ and let $U_{1},\ldots,U_{N}$ be independent
uniforms on $[0,1]$.  For $N\ge1$ and every real $x$,
\begin{equation}
 \Prob\Bigl(\sum_{j=1}^{N}w_{j}U_{j}\le x\Bigr)
 =\frac{1}{N!\prod_{j=1}^{N}w_{j}}
  \sum_{\delta\in\SetOf{0,1}^{N}}(-1)^{\AbsoluteValue{\delta}}
  \PositivePart{x-\sum_{j=1}^{N}w_{j}\delta_{j}}^{N},
 \label{rf:eq:master-box}
\end{equation}
where $\AbsoluteValue{\delta}=\sum_{j}\delta_{j}$ and $s_{+}=\max(s,0)$.
\end{lemma}

\begin{proof}
After the substitution $y_{j}=w_{j}U_{j}$, the probability is the volume of
the part of the box $\prod_{j}[0,w_{j}]$ on which $\sum_{j}y_{j}\le x$,
divided by $\prod_{j}w_{j}$.  The unrestricted nonnegative simplex
$\SetOf{y\ge0:\sum_{j}y_{j}\le x}$ has volume $x_{+}^{N}/N!$, by scaling the
unit simplex, whose volume is $1/N!$ by $N$ successive integrations.  Exclude
the conditions $y_{j}>w_{j}$ by inclusion--exclusion: for a fixed set
$S=\SetOf{j:\delta_{j}=1}$ of violated constraints, translate $y_{j}$ by
$w_{j}$ for $j\in S$; the remaining region is again a nonnegative simplex,
with threshold $x-\sum_{j\in S}w_{j}$.  Summing its volume with the sign
$(-1)^{\AbsoluteValue{S}}$ gives \cref{rf:eq:master-box}.  Boundaries have
zero volume, so the weak inequalities cause no ambiguity.
\end{proof}

For $w_{j}=2^{-j}$ the subset sums are $h/2^{N}$ with $0\le h<2^{N}$: reading
the digits of $h$ from the most significant end, $h=\sum_{j=1}^{N}\delta_{j}2^{N-j}$,
and the sign $(-1)^{\AbsoluteValue{\delta}}$ is the Thue--Morse sign
$\tm{h}$.  The product of the weights is $2^{-\Tr{N}}$.  Hence the
distribution function $H_{N}$ of the uncentred prefix $X_{N}$ is
\begin{equation}
 H_{N}(x)=\frac{2^{\Tr{N}}}{N!}\sum_{h=0}^{2^{N}-1}\tm{h}
  \PositivePart{x-\frac{h}{2^{N}}}^{N}.
 \label{rf:eq:master-uncentred-spline}
\end{equation}
This is a finite formula for an approximant, not yet a finite formula for
$\Fb$.

\subsection{Why the tail becomes polynomial at a dyadic point}
\label{rf:sec:master-identity}

The independent-copy decomposition $X\EqualInLaw X_{n}+2^{-n}X'$, with $X'$
a copy of $X$ independent of $X_{n}$, gives
$\Fb(x)=\E H_{n}(x-2^{-n}X')$ for every real $x$.  At $x=m/2^{n}$ the active
set in \cref{rf:eq:master-uncentred-spline} is particularly simple: the
argument $m-h-X'$ is nonnegative for every $X'\in[0,1]$ when $h<m$ and
nonpositive when $h\ge m$.  The support $0\le X'\le1$ of the tail is exactly
what makes this possible.

\begin{theorem}[Dyadic master identity]
\label{rf:thm:master}
For $n\ge1$ and $0\le m\le2^{n}$,
\begin{align}
 \Fb\!\left(\frac{m}{2^{n}}\right)
 &=2^{-\binom{n}{2}}\sum_{h=0}^{m-1}\tm{h}
   \sum_{l=0}^{n}A_{l}\,\frac{(m-h-1)^{n-l}}{(n-l)!}
 \label{rf:eq:master-raw}\\
 &=2^{-\Tr{n}}\sum_{h=0}^{m-1}\tm{h}
   \sum_{k=0}^{\Floor{n/2}}a_{k}\,\frac{(2m-2h-1)^{n-2k}}{(n-2k)!}
 \label{rf:eq:master-centred}\\
 &=\frac{2^{-\Tr{n}}}{n!}\sum_{h=0}^{m-1}\tm{h}
   \sum_{k=0}^{\Floor{n/2}}\binom{n}{2k}c_{k}\,(2m-2h-1)^{n-2k}.
 \label{rf:eq:master-classical}
\end{align}
The centred form also holds for $n=0$ with the empty sum and $a_{0}=1$,
giving $\Fb(0)=0$ and $\Fb(1)=1$.
\end{theorem}

\begin{proof}
Substitute $x=m/2^{n}$ in $\Fb(x)=\E H_{n}(x-2^{-n}X')$ and use
\cref{rf:eq:master-uncentred-spline} at $N=n$:
\[
 \Fb\!\left(\frac{m}{2^{n}}\right)
 =\frac{2^{\Tr{n}}}{n!}\sum_{h=0}^{2^{n}-1}\tm{h}\,
  \E\PositivePart{\frac{m-h-X'}{2^{n}}}^{n}
 =\frac{2^{\Tr{n}-n^{2}}}{n!}\sum_{h=0}^{m-1}\tm{h}\,
  \E(m-h-X')^{n},
\]
where the inactive terms $h\ge m$ were dropped and the positive part removed
for $h<m$; note $\Tr{n}-n^{2}=-\binom{n}{2}$.  By
\cref{rf:eq:setup-reflection}, $X'\EqualInLaw1-X'$, so the expectation
equals $\E(m-h-1+X)^{n}$; expanding the binomial and using
$\E X^{l}/l!=A_{l}$ gives \cref{rf:eq:master-raw}.  Next write
$m-h-1+X=\tfrac12(2m-2h-1+Z)$: the extra factor $2^{-n}$ turns
$2^{-\binom{n}{2}}$ into $2^{-\Tr{n}}$, the odd moments of $Z$ vanish, and
$\binom{n}{2k}c_{k}/n!=a_{k}/(n-2k)!$ gives \cref{rf:eq:master-centred} and
\cref{rf:eq:master-classical}.  Symmetry of $Z$ permits either sign in front
of it.  For $n=0$ the centred display reads $\sum_{h<m}\tm{h}$ with
$m\in\SetOf{0,1}$.
\end{proof}

This proof explains why the tail produces only finitely many moments: the
dyadic threshold falls between consecutive head knots, so the truncated power
is an ordinary polynomial throughout the tail's support.  Rationality is
not assumed anywhere.

\begin{remark}[The Taylor-integral proof]
\label{rf:rmk:master-taylor}
The signed extension gives a second proof that also covers numerators
larger than $2^{n}$.  Taylor's formula with integral remainder at the flat
point $0$, of order $n+1$, and the dilation law
\cref{rf:eq:setup-global-derivative} give
\[
 \Fg(x)=\frac{1}{n!}\int_{0}^{x}(x-t)^{n}\Fg^{(n+1)}(t)\,\dd t
 =\frac{2^{-\Tr{n}}}{n!}\int_{0}^{2^{n+1}x}(2^{n+1}x-v)^{n}\Fg(v)\,\dd v
\]
after the substitution $v=2^{n+1}t$.  At $x=a/2^{n}$ the integration
interval is $[0,2a]$, a union of $a$ complete signed bump blocks; on the
$h$th block substitute $v=2h+1+y$ and use \cref{rf:eq:setup-global-block}:
\begin{equation}
 \Fg\!\left(\frac{a}{2^{n}}\right)
 =\frac{2^{-\Tr{n}}}{n!}\sum_{h=0}^{a-1}\tm{h}\int_{-1}^{1}(2a-2h-1-y)^{n}\up(y)\,\dd y
 =2^{-\Tr{n}}\sum_{h=0}^{a-1}\tm{h}\sum_{k=0}^{\Floor{n/2}}a_{k}\frac{(2a-2h-1)^{n-2k}}{(n-2k)!},
 \label{rf:eq:master-global}
\end{equation}
valid for every $a\ge0$ and $n\ge0$.  For $0\le a\le2^{n}$ the left side is
$\Fb(a/2^{n})$ and the display is \cref{rf:eq:master-centred}.
\end{remark}

For later use set
\begin{equation}
 \Phi_{r}(v)\DefinitionEquals\coeff{t}{r}\,\EulerE^{vt}M(t)
 =\sum_{k=0}^{\Floor{r/2}}a_{k}\,\frac{v^{r-2k}}{(r-2k)!},\qquad\Phi_{0}=1,
 \label{rf:eq:master-Phi}
\end{equation}
so that \cref{rf:eq:master-centred} reads
\begin{equation}
 \Fb\!\left(\frac{m}{2^{n}}\right)
 =2^{-\Tr{n}}\sum_{h=0}^{m-1}\tm{h}\,\Phi_{n}(2m-2h-1)
 =2^{-\Tr{n}}\coeff{t}{n}\,M(t)\sum_{h=0}^{m-1}\tm{h}\,\EulerE^{(2m-2h-1)t}.
 \label{rf:eq:master-egf}
\end{equation}
At this stage the moments have merely been isolated.  The next section
replaces every $a_{k}$, and every $A_{l}$, by a finite expression.

% =====================================================================
\section{Eliminating every moment}
\label{rf:sec:moments}

\subsection{Bernoulli numbers without a recurrence}
\label{rf:sec:moments-bernoulli}

Define, for every nonnegative integer $s$,
\begin{equation}
 \beta_{s}\DefinitionEquals\sum_{j=0}^{s}\frac{1}{j+1}
  \sum_{v=0}^{j}(-1)^{v}\binom{j}{v}v^{s},
 \label{rf:eq:moments-bernoulli-finite}
\end{equation}
with $0^{0}=1$.

\begin{lemma}[Finite Bernoulli identity]
\label{rf:lem:moments-bernoulli}
As formal power series,
$\sum_{s\ge0}\beta_{s}z^{s}/s!=z/(\EulerE^{z}-1)$.  Consequently
$\beta_{s}=B_{s}$ with $B_{1}=-\tfrac12$, and $B_{2r+1}=0$ for $r\ge1$.
\end{lemma}

\begin{proof}
The exponential generating function of the inner finite difference is
$\sum_{s}\frac{z^{s}}{s!}\sum_{v=0}^{j}(-1)^{v}\binom{j}{v}v^{s}=(1-\EulerE^{z})^{j}$,
which has order $j$ in $z$; hence the inner sum vanishes for $j>s$, and the
sum over $j$ in \cref{rf:eq:moments-bernoulli-finite} may be extended to
infinity without changing any coefficient.  With $u=1-\EulerE^{z}$,
\[
 \sum_{j\ge0}\frac{u^{j}}{j+1}=\frac{-\log(1-u)}{u}=\frac{z}{\EulerE^{z}-1}.
\]
All operations are formal and locally finite in each degree.  Finally
$z/(\EulerE^{z}-1)+z/2=\tfrac{z}{2}\coth\tfrac{z}{2}$ is even, which gives
the odd-index vanishing.
\end{proof}

\subsection{Explicit cumulants}
\label{rf:sec:moments-cumulants}

For $r\ge1$ put
\begin{equation}
 \lambda_{r}\DefinitionEquals\frac{2^{2r-1}B_{2r}}{r\,(2r)!\,(4^{r}-1)},\qquad
 \beta^{\ast}_{r}\DefinitionEquals\frac{\lambda_{r}}{4^{r}}
 =\frac{B_{2r}}{2r\,(2r)!\,(4^{r}-1)}.
 \label{rf:eq:moments-lambda}
\end{equation}
The Bernoulli numbers on the right can always be replaced by
\cref{rf:eq:moments-bernoulli-finite}.

\begin{proposition}[Explicit cumulants]
\label{rf:prop:moments-logM}
Near zero, and also formally,
\begin{equation}
 \log M(t)=\sum_{r\ge1}\lambda_{r}t^{2r},\qquad
 \log M_{X}(z)=\frac{z}{2}+\sum_{r\ge1}\beta^{\ast}_{r}z^{2r}.
 \label{rf:eq:moments-logM}
\end{equation}
Thus the even cumulant of $Z$ of order $2r$ is $(2r)!\lambda_{r}$, every odd
cumulant of $Z$ vanishes, and the cumulants of $X$ beyond the mean $\tfrac12$
are $(2r)!\beta^{\ast}_{r}$.  The first values are
\begin{equation}
 \lambda_{1}=\frac{1}{18},\quad\lambda_{2}=-\frac{1}{2700},\quad
 \lambda_{3}=\frac{1}{178605},\quad\lambda_{4}=-\frac{1}{9639000},\quad
 \lambda_{5}=\frac{1}{478533825}.
 \label{rf:eq:moments-lambda-values}
\end{equation}
\end{proposition}

\begin{proof}
Differentiate $\log H(t)$:
\[
 \frac{\dd}{\dd t}\log H(t)=\coth t-\frac1t
 =1+\frac{2}{\EulerE^{2t}-1}-\frac1t
 =\frac1t\Bigl(\frac{2t}{\EulerE^{2t}-1}+t-1\Bigr)
 =\sum_{r\ge1}\frac{2^{2r}B_{2r}}{(2r)!}t^{2r-1},
\]
by \cref{rf:lem:moments-bernoulli} at $z=2t$: the constant term cancels
because $B_{1}=-\tfrac12$, and the odd Bernoulli numbers vanish.
Integrating with zero constant term gives
$\log H(t)=\sum_{r\ge1}\frac{2^{2r-1}B_{2r}}{r(2r)!}t^{2r}$.  Summing this
at $t/2^{j}$ over $j\ge1$ multiplies the coefficient of $t^{2r}$ by
$\sum_{j\ge1}4^{-rj}=1/(4^{r}-1)$; local uniform convergence of the product
justifies the analytic argument, and coefficientwise geometric summation
gives the same formal result.  The raw statement follows from
$M_{X}(z)=\EulerE^{z/2}M(z/2)$ and $\lambda_{r}4^{-r}=\beta^{\ast}_{r}$.
\end{proof}

Bernoulli numbers need not be accepted even as finite sums: the cumulant
coefficients also have a formula using only factorials.

\begin{lemma}[Bernoulli-free cumulant formula]
\label{rf:lem:moments-lambda-nob}
For every $r\ge1$,
\begin{equation}
 \boxed{\;
 \lambda_{r}=\frac{1}{4^{r}-1}\sum_{\ell=1}^{r}\frac{(-1)^{\ell-1}}{\ell}
 \sum_{\substack{s_{1},\ldots,s_{\ell}\ge1\\ s_{1}+\cdots+s_{\ell}=r}}
 \prod_{i=1}^{\ell}\frac{1}{(2s_{i}+1)!}\,.\;}
 \label{rf:eq:moments-lambda-nob}
\end{equation}
\end{lemma}

\begin{proof}
Write $H(t)=1+u(t)$ with $u(t)=\sum_{s\ge1}t^{2s}/(2s+1)!$, a formal series
without constant term.  Then
$\log H(t)=\sum_{\ell\ge1}\frac{(-1)^{\ell-1}}{\ell}u(t)^{\ell}$, and the
coefficient of $t^{2r}$ in $u^{\ell}$ is the sum over ordered
$\ell$-tuples of positive integers with sum $r$ of
$\prod_{i}1/(2s_{i}+1)!$; only $\ell\le r$ contributes.  Hence
$\coeff{t}{2r}\log H(t)$ is the double sum in
\cref{rf:eq:moments-lambda-nob} without the prefactor.  Summing over the
scales $t/2^{j}$, $j\ge1$, as in the previous proof multiplies it by
$1/(4^{r}-1)$.  Comparison with \cref{rf:eq:moments-logM} gives the claim
(and, incidentally, a factorial-only formula for $2^{2r-1}B_{2r}/(r(2r)!)$).
\end{proof}

\subsection{Multiplicity-profile coefficients}
\label{rf:sec:moments-partitions}

\begin{theorem}[Partition formulas for the moment coefficients]
\label{rf:thm:moments-partition}
For every $k,m\ge0$,
\begin{align}
 a_{k}&=\sum_{(v_{1},\ldots,v_{k})\in\Prof{k}}
        \prod_{r=1}^{k}\frac{\lambda_{r}^{v_{r}}}{v_{r}!},
 \label{rf:eq:moments-a-partition}\\
 A_{m}&=\sum_{\substack{v_{0},v_{1},\ldots,v_{\Floor{m/2}}\ge0\\
                        v_{0}+2\sum_{r\ge1}rv_{r}=m}}
        \frac{2^{-v_{0}}}{v_{0}!}
        \prod_{r=1}^{\Floor{m/2}}\frac{(\beta^{\ast}_{r})^{v_{r}}}{v_{r}!}.
 \label{rf:eq:moments-A-partition}
\end{align}
At index zero each sum consists of the empty or zero vector and has value
one.
\end{theorem}

\begin{proof}
Exponentiating \cref{rf:eq:moments-logM} gives
$M(t)=\prod_{r\ge1}\exp(\lambda_{r}t^{2r})$.  To extract the coefficient of
$t^{2k}$ only the factors with $r\le k$ and the powers $v_{r}\le k/r$ can
contribute; expanding these finitely many exponentials to the relevant
orders and choosing $v_{r}$ copies of $\lambda_{r}t^{2r}$ gives exactly
\cref{rf:eq:moments-a-partition}.  The second identity follows in the same
way from $M_{X}(z)=\EulerE^{z/2}\prod_{r}\exp(\beta^{\ast}_{r}z^{2r})$, the
factor $\EulerE^{z/2}$ contributing $v_{0}$ copies of $z/2$.  The same
calculation is valid coefficientwise in formal power series, so no analytic
evaluation is required.  These are the multiplicity-vector versions of the
complete Bell-polynomial identities, written out so that no convention for
Bell polynomials is needed.
\end{proof}

\begin{theorem}[Fully expanded Bernoulli partition value]
\label{rf:thm:partition-value}
Let $n\ge0$, $0\le m\le2^{n}$, and $d=\Floor{n/2}$.  Then
\begin{equation}
 \boxed{\;
 \Fb\!\left(\frac{m}{2^{n}}\right)
 =2^{-\Tr{n}}\sum_{h=0}^{m-1}\tm{h}
  \sum_{\substack{v_{0},v_{1},\ldots,v_{d}\ge0\\ v_{0}+2\sum_{r=1}^{d}rv_{r}=n}}
  \frac{(2m-2h-1)^{v_{0}}}{v_{0}!}\prod_{r=1}^{d}\frac{\lambda_{r}^{v_{r}}}{v_{r}!}\,,\;}
 \label{rf:eq:partition-value}
\end{equation}
where every $\lambda_{r}$ is the rational expression
\cref{rf:eq:moments-lambda}, with $B_{2r}$ given by
\cref{rf:eq:moments-bernoulli-finite} or $\lambda_{r}$ by
\cref{rf:eq:moments-lambda-nob}.
\end{theorem}

\begin{proof}
Substitute \cref{rf:eq:moments-a-partition} into
\cref{rf:eq:master-centred} and set $v_{0}=n-2k$; conversely every vector in
\cref{rf:eq:partition-value} determines $k=\sum_{r}rv_{r}$, so the reindexing
is a bijection.  Equivalently, expand $\EulerE^{vt}\exp(\sum_{r}\lambda_{r}t^{2r})$
in \cref{rf:eq:master-egf}.
\end{proof}

\subsection{Positive ordered compositions: no Bernoulli numbers}
\label{rf:sec:moments-compositions}

\begin{theorem}[Positive composition formulas]
\label{rf:thm:moments-composition}
For every $k,m\ge0$, with the suffix sums \cref{rf:eq:intro-suffix},
\begin{align}
 a_{k}&=\sum_{\mathbf r\in\Comp{k}}\prod_{i=1}^{\ell(\mathbf r)}
       \frac{1}{(2r_{i}+1)!\,(4^{R_{i}}-1)},
 \label{rf:eq:moments-a-composition}\\
 A_{m}&=\sum_{\mathbf r\in\Comp{m}}\prod_{i=1}^{\ell(\mathbf r)}
       \frac{1}{(r_{i}+1)!\,(2^{R_{i}}-1)}.
 \label{rf:eq:moments-A-composition}
\end{align}
Every summand is positive.  Together with $a_{0}=A_{0}=1$ this supplies
every coefficient in \cref{rf:thm:master} without any special-number
sequence.
\end{theorem}

\begin{proof}
Expand the product \cref{rf:eq:setup-M} using
$H(t/2^{j})=\sum_{r\ge0}4^{-jr}t^{2r}/(2r+1)!$.  In a contribution to degree
$2k$, let $j_{1}<\cdots<j_{\ell}$ be the distinct occupied scales and
$r_{i}\ge1$ the exponent chosen at scale $j_{i}$, so
$(r_{1},\ldots,r_{\ell})\in\Comp{k}$.  Its weight is
$\prod_{i}4^{-j_{i}r_{i}}/(2r_{i}+1)!$.  All coefficients are nonnegative,
so the terms may be regrouped by the composition (monotone convergence, or
first do the argument for the finite product $M_{N}$ and let $N\to\infty$,
the even coefficients increasing to those of $M$).  Put $g_{1}=j_{1}$ and
$g_{i}=j_{i}-j_{i-1}$ for $i>1$: the gaps are positive integers in bijection
with the increasing tuples, and $\sum_{i}j_{i}r_{i}=\sum_{i}g_{i}R_{i}$,
since $j_{i}=g_{1}+\cdots+g_{i}$ and each $g_{i}$ multiplies exactly the
exponents $r_{i},\ldots,r_{\ell}$.  Hence the sum over occupied scales
factors into finite geometric sums,
\begin{equation}
 \sum_{1\le j_{1}<\cdots<j_{\ell}}4^{-\sum_{i}j_{i}r_{i}}
 =\prod_{i=1}^{\ell}\sum_{g_{i}\ge1}4^{-g_{i}R_{i}}
 =\prod_{i=1}^{\ell}\frac{1}{4^{R_{i}}-1},
 \label{rf:eq:moments-gaps}
\end{equation}
which proves \cref{rf:eq:moments-a-composition}.  For the raw coefficients
expand $M_{X}(z)=\prod_{j\ge1}\E\,\EulerE^{2^{-j}zU_{j}}$ with
$\E U^{r}/r!=1/(r+1)!$; the gap calculation is identical with $2$ in place
of $4$, giving \cref{rf:eq:moments-A-composition}.
\end{proof}

\begin{example}[First coefficients]
\label{rf:ex:moments-first}
The compositions $(1)$, and $(2)$, $(1,1)$, give
\[
 a_{1}=\frac{1}{3!\,(4-1)}=\frac{1}{18},\qquad
 a_{2}=\frac{1}{5!\,(16-1)}+\frac{1}{3!\,3!\,(16-1)(4-1)}=\frac{19}{16200},
\]
so $\E Z^{2}=\tfrac19$ and $\E Z^{4}=\tfrac{19}{675}$; the four
compositions of $3$ give $a_{3}=583/42865200$, and $a_{4}=132809/1311675120000$.
The same numbers result from \cref{rf:eq:moments-a-partition}:
$a_{2}=\lambda_{2}+\lambda_{1}^{2}/2$ and
$a_{3}=\lambda_{3}+\lambda_{1}\lambda_{2}+\lambda_{1}^{3}/6$.  Raw:
$A_{1}=\tfrac12$, $A_{2}=\tfrac5{36}$, $A_{3}=\tfrac1{36}$,
$A_{4}=\tfrac{143}{32400}$.  Multiplying normalized coefficients by
$(2k)!$ is essential; omitting the factor would change the dyadic values.
\end{example}

\begin{theorem}[Bernoulli-free closed dyadic sum]
\label{rf:thm:composition-value}
For $n\ge0$, $0\le m\le2^{n}$, $d=\Floor{n/2}$,
\begin{equation}
 \boxed{\begin{aligned}
 \Fb\!\left(\frac{m}{2^{n}}\right)={}&2^{-\Tr{n}}\sum_{h=0}^{m-1}\tm{h}
   \sum_{k=0}^{d}\frac{(2m-2h-1)^{n-2k}}{(n-2k)!}\\[-0.2em]
 &\times\sum_{\ell=0}^{k}\sum_{\substack{r_{1},\ldots,r_{\ell}\ge1\\ r_{1}+\cdots+r_{\ell}=k}}
   \prod_{i=1}^{\ell}\frac{1}{(2r_{i}+1)!\,\bigl(4^{r_{i}+\cdots+r_{\ell}}-1\bigr)}\,.
 \end{aligned}}
 \label{rf:eq:composition-value}
\end{equation}
The $\ell=0$ inner sum is one for $k=0$ and zero otherwise.  Every quantity
is a factorial, an integer power, a Thue--Morse sign, or a bounded sum or
product of these.
\end{theorem}

\begin{proof}
Insert \cref{rf:eq:moments-a-composition} into \cref{rf:eq:master-centred}.
\end{proof}

\begin{remark}[Relation to the triangular recurrence]
\label{rf:rmk:moments-recurrence}
Comparing coefficients of $t^{2k}$ in $M(2t)=H(t)M(t)$ gives
\begin{equation}
 (4^{k}-1)\,a_{k}=\sum_{j=0}^{k-1}\frac{a_{j}}{(2(k-j)+1)!}\qquad(k\ge1).
 \label{rf:eq:moments-recurrence}
\end{equation}
Unrolling this recurrence produces a path interpretation with the same
suffix-composition weights as \cref{rf:eq:moments-a-composition}, so the
composition formula is the closed form of the recurrence, not a recurrence
in disguise.  Equation \cref{rf:eq:moments-recurrence} is retained as an
independent check and as the source of the matrix in
\cref{rf:sec:determinant}; it is not part of any evaluation prescription.
\end{remark}

The composition sum has $2^{k-1}$ terms for $k\ge1$, all positive; the
partition sum has one term per integer partition of $k$, at the price of
alternating signs among the $\lambda_{r}$.  Both are exact rational
formulas, and neither requires $a_{0},\ldots,a_{k-1}$ before $a_{k}$.

\subsection{Finite reciprocal coefficients}
\label{rf:sec:moments-reciprocal}

The scale-polynomial results of \cref{rf:sec:splines} use
\begin{equation}
 \frac{1}{M(t)}=\sum_{k\ge0}e_{k}t^{2k},\qquad e_{0}=1,
 \label{rf:eq:moments-inverse-M}
\end{equation}
as a formal series or near zero.  Exactly the calculation of
\cref{rf:thm:moments-partition} applied to $\exp(-\sum_{r}\lambda_{r}t^{2r})$
gives the recurrence-free formula
\begin{equation}
 e_{k}=\sum_{(v_{1},\ldots,v_{k})\in\Prof{k}}
       \prod_{r=1}^{k}\frac{(-\lambda_{r})^{v_{r}}}{v_{r}!},
 \label{rf:eq:moments-e-partition}
\end{equation}
with first values
\begin{equation}
 e_{1}=-\frac{1}{18},\quad e_{2}=\frac{31}{16200},\quad
 e_{3}=-\frac{2347}{42865200},\quad e_{4}=\frac{1904369}{1311675120000}.
 \label{rf:eq:moments-e-values}
\end{equation}
Alternatively, the $\lambda_{r}$ can be avoided by the finite geometric
expansion of $1/(1+(M-1))$:
\begin{equation}
 e_{k}=\sum_{\ell=1}^{k}(-1)^{\ell}
 \sum_{\substack{r_{1},\ldots,r_{\ell}\ge1\\ r_{1}+\cdots+r_{\ell}=k}}
 \prod_{i=1}^{\ell}a_{r_{i}}\qquad(k\ge1),
 \label{rf:eq:moments-e-compositions}
\end{equation}
followed by \cref{rf:eq:moments-a-composition}.  The signs $(-1)^{k}e_{k}>0$
are proved in \cref{rf:lem:splines-e-sign}.

\subsection{Reciprocal powers of two}
\label{rf:sec:moments-reciprocal-powers}

At $m=1$ the raw form \cref{rf:eq:master-raw} collapses to a single
coefficient:
\begin{equation}
 \boxed{\;\Fb(2^{-n})=2^{-\binom{n}{2}}A_{n}\;}
 \qquad(n\ge0),
 \label{rf:eq:moments-inverse-power}
\end{equation}
so either \cref{rf:eq:moments-A-partition} or
\cref{rf:eq:moments-A-composition} is a direct finite formula for these
values, for example $\Fb(\tfrac1{16})=2^{-6}A_{4}=143/2073600$.  In centred
form, \cref{rf:eq:master-centred} at $m=1$ gives the all-positive expression
\begin{equation}
 \Fb(2^{-n})=2^{-\Tr{n}}\sum_{k=0}^{\Floor{n/2}}\frac{1}{(n-2k)!}
 \sum_{\mathbf r\in\Comp{k}}\prod_{i=1}^{\ell(\mathbf r)}
 \frac{1}{(2r_{i}+1)!\,(4^{R_{i}}-1)}.
 \label{rf:eq:moments-inverse-power-positive}
\end{equation}
For odd depths there is a further simplification through the bump's half
moments:
\begin{equation}
 \Fb\!\left(2^{-(2j+1)}\right)=\frac{c_{j}}{2\,(2j)!\,2^{\binom{2j+1}{2}}}
 =\frac{a_{j}}{2^{\binom{2j+1}{2}+1}}\qquad(j\ge0).
 \label{rf:eq:moments-odd-depth}
\end{equation}
To prove it without a moment recurrence, note that
$\up(t)=\Fb(1-t)=\Prob(X\ge t)$ on $[0,1]$ by
\cref{rf:eq:setup-fold,rf:eq:setup-reflection}.  Layer-cake integration
gives $\E X^{2j+1}=(2j+1)\int_{0}^{1}t^{2j}\Prob(X\ge t)\,\dd t
=(2j+1)\int_{0}^{1}t^{2j}\up(t)\,\dd t=(2j+1)\,c_{j}/2$, by evenness of
$\up$ and $\int_{-1}^{1}t^{2j}\up=c_{j}$; hence
$A_{2j+1}=c_{j}/(2\,(2j)!)$, and \cref{rf:eq:moments-inverse-power} gives
the claim.  For example $\Fb(\tfrac18)=a_{1}/2^{4}=1/288$ and
$\Fb(\tfrac1{32})=a_{2}/2^{11}=19/33177600$.

% =====================================================================
\section{Binary compression of the numerator}
\label{rf:sec:binary}

The master formula sums over $m$ consecutive integers.  Its special signs
permit an exact compression into binary blocks, in two different ways, and
both are useful independently of which coefficient formula is chosen.

\subsection{One outer summand per nonzero binary digit}
\label{rf:sec:binary-block}

Write the binary expansion and the digit data
\begin{equation}
 m=\sum_{j\ge0}b_{j}2^{j},\quad b_{j}\in\SetOf{0,1},\qquad
 J(m)\DefinitionEquals\SetOf{j:b_{j}=1},
 \label{rf:eq:binary-digits}
\end{equation}
and for $j\in J(m)$
\begin{equation}
 p_{j}\DefinitionEquals2^{j+1}\Floor{\frac{m}{2^{j+1}}},\qquad
 r_{j}\DefinitionEquals m\bmod2^{j},\qquad
 K_{j}\DefinitionEquals2^{j}+2r_{j},\qquad
 \eta_{j}\DefinitionEquals\tm{\Floor{m/2^{j+1}}}.
 \label{rf:eq:binary-parameters}
\end{equation}
Thus $p_{j}$ is the part of $m$ strictly above the $j$th bit, $r_{j}$ the
part strictly below it, and $m=p_{j}+2^{j}+r_{j}$; these are digit
operations, not recursive values of $\Fb$.

\begin{theorem}[Binary block formula]
\label{rf:thm:binary-block}
For every $a\ge0$ and $n\ge0$, including numerators larger than $2^{n}$,
\begin{equation}
 \boxed{\;
 \Fg\!\left(\frac{a}{2^{n}}\right)
 =2^{-\Tr{n}}\sum_{\substack{j\in J(a)\\ j\le n}}\eta_{j}\,2^{\Tr{j}}
  \sum_{k=0}^{\Floor{(n-j)/2}}
  \frac{a_{k}\,4^{jk}\,K_{j}^{\,n-j-2k}}{(n-j-2k)!}\,,\;}
 \label{rf:eq:binary-block}
\end{equation}
equivalently, with $\Phi_{r}$ from \cref{rf:eq:master-Phi} and
$A_{j}\DefinitionEquals K_{j}/2^{j}=1+2^{1-j}r_{j}$,
\begin{equation}
 \Fg\!\left(\frac{a}{2^{n}}\right)
 =\sum_{\substack{j\in J(a)\\ j\le n}}\eta_{j}\,2^{-\Tr{n-j}}\,\Phi_{n-j}(A_{j}).
 \label{rf:eq:binary-block-Phi}
\end{equation}
For $0\le a\le2^{n}$ the left side is $\Fb(a/2^{n})$.  The coefficients
$a_{k}$ may be supplied by \cref{rf:eq:moments-a-composition} or
\cref{rf:eq:moments-a-partition}; the outer sum has at most
$\min(n+1,\BinaryDigitSum(a))$ nonzero terms, independently of the size of
$a$.
\end{theorem}

\begin{proof}
The interval of integers $[0,a)$ is the disjoint union of the blocks
$[p_{j},p_{j}+2^{j})$, $j\in J(a)$.  Adding an integer $h<2^{j}$ to $p_{j}$
causes no binary carry, so $\tm{p_{j}+h}=\tm{p_{j}}\tm{h}=\eta_{j}\tm{h}$.
The finite Thue--Morse product identity
\begin{equation}
 \sum_{h=0}^{2^{j}-1}\tm{h}\,z^{h}=\prod_{\ell=0}^{j-1}\bigl(1-z^{2^{\ell}}\bigr)
 \label{rf:eq:binary-tm-product}
\end{equation}
holds because selecting a subset of the factors selects the nonzero digits
of $h$.  Therefore
\begin{equation}
 \sum_{h=0}^{a-1}\tm{h}\,z^{h}
 =\sum_{j\in J(a)}\eta_{j}\,z^{p_{j}}\prod_{\ell=0}^{j-1}\bigl(1-z^{2^{\ell}}\bigr).
 \label{rf:eq:binary-prefix-product}
\end{equation}
By \cref{rf:eq:master-global}, the desired value is
$2^{-\Tr{n}}\coeff{t}{n}\,\EulerE^{(2a-1)t}M(t)\sum_{h<a}\tm{h}\EulerE^{-2ht}$.
For the block of length $2^{j}$, each factor of the product at
$z=\EulerE^{-2t}$ is
$1-\EulerE^{-2^{\ell+1}t}=2^{\ell+1}t\,\EulerE^{-2^{\ell}t}H(2^{\ell}t)$, so
\begin{equation}
 \prod_{\ell=0}^{j-1}\bigl(1-\EulerE^{-2^{\ell+1}t}\bigr)
 =2^{\Tr{j}}\,t^{j}\,\EulerE^{-(2^{j}-1)t}\prod_{\ell=0}^{j-1}H(2^{\ell}t),
 \label{rf:eq:binary-H-product}
\end{equation}
the powers of two multiplying to $2^{j+j(j-1)/2}=2^{\Tr{j}}$.  Iterating
$M(2t)=H(t)M(t)$ gives the exact cancellation
$M(t)\prod_{\ell=0}^{j-1}H(2^{\ell}t)=M(2^{j}t)$.  The block therefore
contributes
\[
 \eta_{j}\,2^{-\Tr{n}+\Tr{j}}\coeff{t}{n-j}\,\EulerE^{(2a-2p_{j}-2^{j})t}M(2^{j}t),
\]
and $2a-2p_{j}-2^{j}=2^{j}+2r_{j}=K_{j}$.  Expanding
$M(2^{j}t)=\sum_{k}a_{k}4^{jk}t^{2k}$ and extracting $t^{n-j}$ proves
\cref{rf:eq:binary-block}; a block with $j>n$ carries the factor $t^{j}$ and
contributes nothing, which is why the high bits enter only through the signs
$\eta_{j}$.  Substituting $s=2^{j}t$ changes the coefficient extraction by
the factor $2^{j(n-j)}$, turns the exponent into $A_{j}s$, and
$-\Tr{n}+\Tr{j}+j(n-j)=-\Tr{n-j}$ gives \cref{rf:eq:binary-block-Phi}.
\end{proof}

For $a=0$ the empty sum gives zero; for $a=2^{n}$ the single block gives
$\Phi_{0}(1)=1$; for $a=1$ only the lowest bit occurs and
\cref{rf:eq:binary-block-Phi} reduces to $\Fb(2^{-n})=2^{-\Tr{n}}\Phi_{n}(1)$,
which is \cref{rf:eq:moments-inverse-power-positive}.  Combining
\cref{rf:eq:binary-block} with \cref{rf:eq:moments-a-composition} gives one
literal expression with no unresolved coefficient symbol,
\begin{equation}
 \boxed{\begin{aligned}
 \Fg\!\left(\frac{a}{2^{n}}\right)
 ={}&\sum_{\substack{j\in J(a)\\ j\le n}}\eta_{j}\,2^{\Tr{j}-\Tr{n}}
   \sum_{k=0}^{\Floor{(n-j)/2}}\frac{4^{jk}K_{j}^{\,n-j-2k}}{(n-j-2k)!}\\[-0.1em]
 &\times\sum_{\ell=0}^{k}\sum_{\substack{r_{1},\ldots,r_{\ell}\ge1\\ r_{1}+\cdots+r_{\ell}=k}}
   \prod_{i=1}^{\ell}\frac{1}{(2r_{i}+1)!\,\bigl(4^{r_{i}+\cdots+r_{\ell}}-1\bigr)}\,,
 \end{aligned}}
 \label{rf:eq:binary-fully-expanded}
\end{equation}
which, with the digit definitions \cref{rf:eq:binary-parameters}, contains
no recursively defined constant at all.  After the universal coefficients
have been computed there are at most $\sum_{j=0}^{n}(\Floor{(n-j)/2}+1)
=(d+1)(n-d+1)$ argument-dependent monomials --- a count of rational
operations, not of bit complexity.

\begin{example}[Three nonzero bits]
\label{rf:ex:binary-21}
For $x=21/64$, $21=(10101)_{2}$ and $J(21)=\SetOf{0,2,4}$.  The block data
give
\begin{align*}
 \Fb\!\left(\frac{21}{64}\right)
 &=2^{-3}\Phi_{2}\!\left(\frac{13}{8}\right)
 -2^{-10}\Phi_{4}\!\left(\frac{3}{2}\right)
 +2^{-21}\Phi_{6}(1)\\
 &=\frac{1585}{9216}-\frac{71179}{265420800}+\frac{1153}{561842749440}
 =\frac{482385079229}{2809213747200}.
\end{align*}
For $a=3$, $n=3$ the two blocks $(\eta_{0},K_{0})=(-1,1)$ and
$(\eta_{1},K_{1})=(1,4)$ give
$\Fb(\tfrac38)=\tfrac1{64}\bigl[-(\tfrac16+\tfrac1{18})+2(\tfrac{16}{2}+\tfrac4{18})\bigr]
=\tfrac{73}{288}$.  Without a prefix sum of length $173$,
\begin{equation}
 \Fb\!\left(\frac{173}{4096}\right)
 =\frac{1490588447599319294372326859976672988757}
       {292214732887898713986916575925267070976000000}
 \approx5.101003747717089\cdot10^{-6},
 \label{rf:eq:binary-large}
\end{equation}
where only the five set bits of $173=128+32+8+4+1$ contribute.
\end{example}

\subsection{The binary telescope}
\label{rf:sec:binary-telescope}

A second compression uses raw coefficients and the signed extension's
antiperiodicity rather than the block product; it is a finite telescope of
the exact Taylor reduction used in the classical arithmetic theory
\cite{arias2018}.

\begin{lemma}[One binary step]
\label{rf:lem:binary-step}
For an integer $b\ge1$ and $0\le t\le2^{-b}$,
\begin{equation}
 \Fb(2^{-b}+t)+\Fb(t)=2^{-\binom{b}{2}}\sum_{k=0}^{b}A_{b-k}\frac{(2^{b}t)^{k}}{k!}.
 \label{rf:eq:binary-step}
\end{equation}
\end{lemma}

\begin{proof}
On $[0,2]$, \cref{rf:eq:setup-global-block} and $\tm{1}=-1$ give
$\Fg(z+2)=-\Fg(z)$.  By \cref{rf:eq:setup-global-derivative} the
$(b+1)$st derivative of the left side of \cref{rf:eq:binary-step} is
$2^{\Tr{b+1}}\bigl[\Fg(2+2^{b+1}t)+\Fg(2^{b+1}t)\bigr]=0$ for
$0\le2^{b+1}t\le2$, so the left side is a polynomial of degree at most $b$
on the stated interval.  At $t=0$ every derivative of $\Fb(t)$ vanishes,
while for $0\le k\le b$
\[
 \Fb^{(k)}(2^{-b})=2^{\Tr{k}}\Fb\bigl(2^{-(b-k)}\bigr)
 =2^{\Tr{k}-\binom{b-k}{2}}A_{b-k}=2^{bk-\binom{b}{2}}A_{b-k}
\]
by \cref{rf:eq:setup-global-derivative} and
\cref{rf:eq:moments-inverse-power}, using
$\Tr{k}-\binom{b-k}{2}=bk-\binom{b}{2}$; the endpoint $k=b$ uses
$\Fb(1)=A_{0}=1$.  The Taylor polynomial of degree $b$ at $t=0$ is therefore
exactly the right side of \cref{rf:eq:binary-step}.
\end{proof}

\begin{theorem}[Explicit binary telescope]
\label{rf:thm:binary-telescope}
Let $0<x<1$ be dyadic, with nonzero binary positions
$x=\sum_{i=1}^{\rho}2^{-b_{i}}$, $1\le b_{1}<\cdots<b_{\rho}$, and the
explicitly known binary suffixes
\begin{equation}
 s_{i}\DefinitionEquals\sum_{j=i+1}^{\rho}2^{b_{i}-b_{j}},\qquad0\le s_{i}<1.
 \label{rf:eq:binary-suffix}
\end{equation}
Then
\begin{equation}
 \boxed{\;
 \Fb(x)=\sum_{i=1}^{\rho}(-1)^{i-1}\,2^{-\binom{b_{i}}{2}}
 \sum_{k=0}^{b_{i}}\frac{s_{i}^{k}}{k!}
 \sum_{\mathbf r\in\Comp{b_{i}-k}}\prod_{j=1}^{\ell(\mathbf r)}
 \frac{1}{(r_{j}+1)!\,(2^{R_{j}}-1)}\,.\;}
 \label{rf:eq:binary-telescope}
\end{equation}
The inner composition sum is $A_{b_{i}-k}$ and may equally be replaced by
\cref{rf:eq:moments-A-partition}.
\end{theorem}

\begin{proof}
Write $x_{i}=\sum_{j=i}^{\rho}2^{-b_{j}}$ and $x_{\rho+1}=0$.  Since
$x_{i}=2^{-b_{i}}+x_{i+1}$ and $2^{b_{i}}x_{i+1}=s_{i}<1$,
\cref{rf:lem:binary-step} gives $\Fb(x_{i})+\Fb(x_{i+1})=P_{i}(s_{i})$ with
$P_{i}$ the polynomial on the right of \cref{rf:eq:binary-step}.  Multiply
by $(-1)^{i-1}$ and sum over $i$: all interior values cancel and
$\Fb(x_{\rho+1})=\Fb(0)=0$.  Substituting
\cref{rf:eq:moments-A-composition} gives the display.  This is a proof by a
finite telescope, not an instruction to evaluate a recursive sequence.
\end{proof}

For example $3/8=2^{-2}+2^{-3}$ gives
$\Fb(\tfrac38)=2^{-1}\bigl(A_{2}+\tfrac12A_{1}+\tfrac18A_{0}\bigr)-2^{-3}A_{3}=\tfrac{73}{288}$.
A point with $\rho$ nonzero bits uses $\sum_{i}(b_{i}+1)\le\rho(n+1)$
polynomial terms once the raw coefficients through index $n$ are available.

\begin{remark}[The two compressions compared]
\label{rf:rmk:binary-compare}
Both formulas have one outer term per set bit, but they are different
identities: the block formula groups the Thue--Morse prefix into complete
blocks, whose sign structure is absorbed by the exact product identity
\cref{rf:eq:binary-H-product}, uses the centred coefficients $a_{k}$ with
the largest index $\Floor{(n-j)/2}$, and holds for the signed extension at
every numerator; the telescope alternates signs down the bits, uses the raw
coefficients $A_{l}$ with index up to $b_{i}$, and is confined to the unit
interval.  Reflection $\Fb(x)=1-\Fb(1-x)$ may reduce the number of useful
bits in either.  Neither count is a bit-complexity estimate: large rational
numerators and the coefficient generation still have costs.
\end{remark}

% =====================================================================
\section{Eliminating moments by finitely many uniform prefixes}
\label{rf:sec:prefix}

The coefficients $a_{k}$ admit a third recurrence-free representation:
instead of expanding the infinite product, interpolate a coefficient of its
finite prefixes $M_{N}$.  The same device turns the whole master identity
into a formula with finitely many uniform coordinates.

\subsection{Two explicit finite-prefix coefficients}
\label{rf:sec:prefix-coefficients}

Write
\begin{equation}
 a_{k,N}\DefinitionEquals\coeff{t}{2k}M_{N}(t),\qquad
 a_{0,0}=1,\quad a_{k,0}=0\ (k>0).
 \label{rf:eq:prefix-coefficient}
\end{equation}
Expanding the $N$ factors gives the finite multinomial formula
\begin{equation}
 a_{k,N}=\sum_{\substack{r_{1},\ldots,r_{N}\ge0\\ r_{1}+\cdots+r_{N}=k}}
 \prod_{i=1}^{N}\frac{4^{-ir_{i}}}{(2r_{i}+1)!},
 \label{rf:eq:prefix-multinomial}
\end{equation}
a finite sum because every $r_{i}\le k$.  There is also a single signed
power sum.  Since $H(t/2^{j})=2^{j-1}t^{-1}(\EulerE^{t/2^{j}}-\EulerE^{-t/2^{j}})$,
expanding the finite product gives, with $L_{N}=1-2^{-N}$,
\begin{equation}
 M_{N}(t)=2^{N(N-1)/2}\,t^{-N}\sum_{h=0}^{2^{N}-1}\tm{h}\,
 \EulerE^{(L_{N}-2^{1-N}h)t}:
 \label{rf:eq:prefix-exponentials}
\end{equation}
the subset sums $\sum_{j\in S}2^{-j}$ run through $h/2^{N}$ with sign
$\tm{h}$ (reversing the order of the $N$ digits does not change the sign),
and the constant is $\prod_{j=1}^{N}2^{j-1}=2^{N(N-1)/2}$.  The apparent
pole is removable, the sum vanishing to order $N$ at $t=0$ by Prouhet's
cancellation \cref{rf:eq:splines-prouhet} below.  Expanding the
exponentials,
\begin{equation}
 \boxed{\;
 a_{k,N}=\frac{2^{N(N-1)/2-N(2k+N)}}{(2k+N)!}
 \sum_{h=0}^{2^{N}-1}\tm{h}\,(2^{N}-1-2h)^{2k+N}\,.\;}
 \label{rf:eq:prefix-power}
\end{equation}
For instance $a_{1,1}=\tfrac1{24}$; the raw second moment $\E Z_{1}^{2}=\tfrac1{12}$
is not $a_{1,1}$, and the normalization by $(2k)!$ remains essential.

\subsection{Exact coefficient extraction at the quarter base}
\label{rf:sec:prefix-extraction}

The geometric weights $w_{d,j}(q)$ are defined and proved in
\cref{rf:sec:extraction-row}; here only the identity
\cref{rf:eq:extraction-filter} is used, which says that
$\sum_{j=0}^{d}w_{d,j}(q)P(zq^{j})=P(0)$ for every polynomial $P$ of degree
at most $d$ and every $z\neq0$.

\begin{theorem}[Finite-prefix moment elimination]
\label{rf:thm:moment-extraction}
For every $k\ge0$ and every integer $N_{0}\ge0$,
\begin{equation}
 \boxed{\;a_{k}=\sum_{j=0}^{k}w_{k,j}\,a_{k,N_{0}+j}\,,\;}
 \label{rf:eq:moment-extraction}
\end{equation}
with $w_{k,j}=w_{k,j}(\tfrac14)$ the explicit rational weights
\cref{rf:eq:main-weights-integer} and each $a_{k,N}$ given by
\cref{rf:eq:prefix-multinomial} or \cref{rf:eq:prefix-power}.  No infinite
moment remains on the right.  A common row of degree $d\ge k$ may replace
the row of degree $k$.
\end{theorem}

\begin{proof}
From $M_{N}(t)=M(t)/M(2^{-N}t)$ and \cref{rf:eq:moments-inverse-M},
\begin{equation}
 a_{k,N}=\sum_{r=0}^{k}a_{k-r}\,e_{r}\,(4^{-N})^{r}.
 \label{rf:eq:prefix-scale-polynomial}
\end{equation}
For fixed $k$ this is a polynomial of degree at most $k$ in $4^{-N}$ with
constant term $a_{k}$.  Apply \cref{rf:eq:extraction-filter} with
$q=\tfrac14$, $z=4^{-N_{0}}$ and the nodes $4^{-(N_{0}+j)}$.  The occurrence
of $a_{k-r}$ in this proof is not an auxiliary recurrence in the final
expression: all sampled coefficients on the right of
\cref{rf:eq:moment-extraction} are supplied by independent finite formulas.
\end{proof}

At $N_{0}=0$ and $k>0$ the $j=0$ term vanishes, so only the prefixes
$1,\ldots,k$ are needed: $a_{1}=\tfrac43a_{1,1}=\tfrac1{18}$.  Inserting
\cref{rf:eq:moment-extraction,rf:eq:prefix-power} into
\cref{rf:eq:master-Phi} and then into \cref{rf:eq:master-centred} or
\cref{rf:eq:binary-block-Phi} produces another family consisting entirely of
finite Thue--Morse sums, in which the interpolation depth is controlled by
the moment index $k$ and the recovered coefficient is reusable at many
dyadic arguments.

\subsection{A moment-free multinomial formula}
\label{rf:sec:prefix-cube}

Applying the same interpolation to the whole master polynomial gives a
formula with finitely many uniform coordinates.  For $N\ge0$ and any $A$,
define
\begin{equation}
 C_{n,N}(A)\DefinitionEquals\frac{1}{n!}\E(A+Z_{N})^{n}
 =\sum_{\substack{r_{0},r_{1},\ldots,r_{N}\ge0\\ r_{0}+2r_{1}+\cdots+2r_{N}=n}}
 \frac{A^{r_{0}}}{r_{0}!}\prod_{j=1}^{N}\frac{4^{-jr_{j}}}{(2r_{j}+1)!},
 \label{rf:eq:prefix-C}
\end{equation}
so $C_{n,0}(A)=A^{n}/n!$ and every summation index is bounded by $n$.

\begin{theorem}[Finite-cube closed form]
\label{rf:thm:finite-cube}
For every $m,n\ge0$ with $0\le m\le2^{n}$ and $d=\Floor{n/2}$,
\begin{equation}
 \boxed{\;
 \Fb\!\left(\frac{m}{2^{n}}\right)
 =2^{-\Tr{n}}\sum_{h=0}^{m-1}\tm{h}\sum_{j=0}^{d}w_{d,j}\,C_{n,j}(2m-2h-1)\,,\;}
 \label{rf:eq:finite-cube}
\end{equation}
and no moment of an infinite distribution occurs in its evaluation.
\end{theorem}

\begin{proof}
With $C_{n}(A)\DefinitionEquals\coeff{t}{n}\EulerE^{At}M(t)=\Phi_{n}(A)$ and
$C_{n,N}(A)=\coeff{t}{n}\EulerE^{At}M(t)/M(2^{-N}t)$, the reciprocal
expansion gives $C_{n,N}(A)=\sum_{r=0}^{d}e_{r}4^{-Nr}C_{n-2r}(A)$, a
polynomial of degree at most $d$ in $4^{-N}$ with constant term
$\Phi_{n}(A)$.  By \cref{rf:eq:extraction-filter} at $z=1$,
$\Phi_{n}(A)=\sum_{j=0}^{d}w_{d,j}C_{n,j}(A)$; substitute this into
\cref{rf:eq:master-egf}.
\end{proof}

The consecutive depths $0,1,\ldots,d$ are convenient, not compulsory: for
distinct nonnegative $N_{0},\ldots,N_{d}$ and $z_{j}=4^{-N_{j}}$, the
Lagrange weights $\prod_{\ell\neq j}(-z_{\ell})/(z_{j}-z_{\ell})$ replace
$w_{d,j}$ (\cref{rf:thm:extraction-arbitrary} in the same form), and for
$N_{j}=N+Lj$ they reduce to $w_{d,j}(4^{-L})$.

% =====================================================================
\section{The exact polynomial hidden in centred finite splines}
\label{rf:sec:splines}

This section explains why a short moment-free formula exists.  It is
important that the cutoff dependence terminates \emph{exactly}, not merely
to a prescribed asymptotic accuracy.

\subsection{The centred finite distribution}
\label{rf:sec:splines-centred}

For $N\ge1$ define
\begin{equation}
 Y_{N}\DefinitionEquals X_{N}+2^{-N-1}=\frac{1+Z_{N}}{2},\qquad
 G_{N}(x)\DefinitionEquals\Prob(Y_{N}\le x).
 \label{rf:eq:splines-centred}
\end{equation}
The added midpoint $2^{-N-1}$ is the mean of the omitted tail; this
centring makes the correction polynomial even in the small scale, a
simplification the uncentred $H_{N}$ does not enjoy
(\cref{rf:sec:halfbase}).  The finite random variable is centred at
$\tfrac12$, like $X$, and $G_{N}$ is a piecewise polynomial of degree $N$
whose knots lie among the points $(h+\tfrac12)/2^{N}$, $0\le h<2^{N}$.  The
box formula gives
\begin{equation}
 G_{N}(x)=\frac{2^{\Tr{N}}}{N!}\sum_{h=0}^{2^{N}-1}\tm{h}
 \PositivePart{x-\frac{h+\frac12}{2^{N}}}^{N}.
 \label{rf:eq:splines-G-positive}
\end{equation}
At $x=m/2^{n}$ with $N\ge n\ge1$, put $A=2^{N-n}m$; exactly the terms
$h<A$ are active, every one a full power, and
\begin{equation}
 \boxed{\;
 G_{N}\!\left(\frac{m}{2^{n}}\right)
 =\frac{1}{2^{\Tr{N}}N!}\sum_{h=0}^{A-1}\tm{h}\,(2A-2h-1)^{N}
 =\frac{\Spow{N}{A}}{2^{\Tr{N}}N!}\,.\;}
 \label{rf:eq:splines-G-integer}
\end{equation}
The threshold is always halfway between two consecutive possible knots.
The sum is empty at $m=0$ and evaluates to one at $m=2^{n}$.  Its negative
terms come only from inclusion--exclusion; $G_{N}$ itself is a distribution
function.

\subsection{A finite expression for every dyadic derivative}
\label{rf:sec:splines-derivatives}

\begin{lemma}[Differentiated master formula]
\label{rf:lem:splines-derivative-master}
For $n\ge1$, $0\le m\le2^{n}$, and $0\le r\le n$,
\begin{equation}
 \Fb^{(r)}\!\left(\frac{m}{2^{n}}\right)
 =2^{(n+1)r-\Tr{n}}\sum_{h=0}^{m-1}\tm{h}\,\Phi_{n-r}(2m-2h-1).
 \label{rf:eq:splines-derivative-master}
\end{equation}
For $r\ge n+1$ the derivative is zero.  In particular every derivative at a
dyadic point is explicitly rational.
\end{lemma}

\begin{proof}
Conditioning on the independent tail gives, for every real $y$,
\[
 \Fb(y)=\frac{2^{\Tr{n}}}{n!}\sum_{h=0}^{2^{n}-1}\tm{h}\,
 \E\PositivePart{y-2^{-n}(h+X')}^{n}.
\]
For $r\le n$ differentiate $r$ times under the expectation; when $r=n$ the
derivative of the truncated power is $n!$ times the indicator of a positive
argument, and $X'$ has no atoms, so the value at the cutoff is immaterial.
At $y=m/2^{n}$ all terms with $h<m$ have a nonnegative argument for every
$X'\in[0,1]$ and the remaining terms vanish.  Replacing $X'$ by
$1-X'=(1-Z)/2$ and using evenness of $Z$, the scale is
$2^{\Tr{n}}\cdot2^{-n(n-r)}\cdot2^{-(n-r)}=2^{(n+1)r-\Tr{n}}$ and the
expectation divided by $(n-r)!$ is $\Phi_{n-r}$.  Higher derivatives vanish
by \cref{rf:lem:setup-dilation}.
\end{proof}

Comparing with \cref{rf:eq:setup-global-derivative} evaluated by the master
identity, this is also the statement $\Fb^{(r)}(x)=2^{\Tr{r}}\Fg(2^{r}x)$
written out; \cref{rf:sec:rvachev-derivatives} returns to it.

\subsection{Finite-jet deconvolution}
\label{rf:sec:splines-deconvolution}

Let $\Dop$ denote differentiation of polynomials in the variable $s$.  The
next lemma is the one analytic step; it is stated generally because it is
used four times (for the CDF, the density, the shifted spline, and the
integer bases).

\begin{lemma}[Local polynomial smoothing is invertible on a finite jet]
\label{rf:lem:splines-local-inverse}
Let $H$ be a piecewise polynomial that agrees with one polynomial $p(s)$ of
degree at most $N$ at the points $x+s$, $\AbsoluteValue{s}\le\eta$, and let
$W$ be a random variable supported on $[-1,1]$ with a bounded density.  Put
$f(t)\DefinitionEquals\E H(t-\eta W)$ and let
$\MomentGeneratingFunction{W}(z)=\E\,\EulerE^{zW}$.  Suppose the derivatives
of $f$ through order $N$ at $x$ are obtained by convolving the
corresponding piecewise derivatives of $H$.  Then, as an identity of
polynomials of degree at most $N$,
\begin{equation}
 \sum_{r=0}^{N}\frac{f^{(r)}(x)}{r!}s^{r}=\MomentGeneratingFunction{W}(-\eta\Dop)\,p(s),
 \label{rf:eq:splines-jet-operator}
\end{equation}
where the operator is truncated at order $N$, and consequently
\begin{equation}
 p(s)=\bigl[\MomentGeneratingFunction{W}(-\eta\Dop)\bigr]^{-1}
 \sum_{r=0}^{N}\frac{f^{(r)}(x)}{r!}s^{r},
 \label{rf:eq:splines-local-inverse}
\end{equation}
the inverse being the finite operator obtained by truncating the formal
reciprocal series at order $N$.  If $W$ is symmetric, only even powers of
$\eta\Dop$ occur.  For the distribution-function splines of this section and
the density splines of \cref{rf:sec:rvachev-density}, the stipulated
derivative identities hold.
\end{lemma}

\begin{proof}
For $\AbsoluteValue{s}\le\eta$ and $\AbsoluteValue{W}\le1$ the point
$x+s-\eta W$ may leave the cell, so the identity is proved at the level of
jets.  Differentiating the convolution as stipulated and evaluating at $x$,
$f^{(r)}(x)=\E H^{(r)}(x-\eta W)=\E p^{(r)}(-\eta W)$ for $0\le r\le N$,
because $x-\eta W$ lies in the cell.  Taylor's formula for the polynomial
$p^{(r)}$ gives $p^{(r)}(-\eta W)=\sum_{i}p^{(r+i)}(0)(-\eta W)^{i}/i!$, and
taking expectations coefficient by coefficient,
$f^{(r)}(x)=\sum_{i}\E[W^{i}](-\eta)^{i}p^{(r+i)}(0)/i!$, which is the
coefficient identity behind \cref{rf:eq:splines-jet-operator}.  On
polynomials of degree at most $N$ one has $\Dop^{N+1}=0$, so
$\MomentGeneratingFunction{W}(-\eta\Dop)$ is a unipotent finite operator
($\MomentGeneratingFunction{W}(0)=1$), invertible with inverse the truncated
reciprocal series; applying it gives \cref{rf:eq:splines-local-inverse}.  No
convergence of a Taylor series of $f$ is assumed anywhere.  For an
$N$-uniform distribution function the derivatives through order $N-1$ are
continuous and the $N$th is bounded and piecewise constant; these are also
its distributional derivatives, so convolution with a bounded density gives
the derivative identities, values at the finitely many knots not affecting
any integral.  The density spline is the same argument one degree lower.
\end{proof}

\subsection{The exact shrinking-cell identity}
\label{rf:sec:splines-cell}

Let
\begin{equation}
 P_{x}(h)\DefinitionEquals\sum_{r=0}^{n}\frac{\Fb^{(r)}(x)}{r!}h^{r},\qquad x=\frac{m}{2^{n}},
 \label{rf:eq:splines-taylor}
\end{equation}
the Taylor polynomial of $\Fb$ at $x$, which terminates by
\cref{rf:eq:setup-jet-termination}; no equality $\Fb(x+h)=P_{x}(h)$ is
asserted.

\begin{theorem}[Exact shrinking-cell formula]
\label{rf:thm:local-cell}
Let $N\ge n\ge1$ and $x=m/2^{n}\in[0,1]$.  For every real $h$ with
$\AbsoluteValue{h}\le2^{-N-1}$,
\begin{equation}
 \boxed{\;
 G_{N}(x+h)=\sum_{k=0}^{\Floor{n/2}}e_{k}\,2^{-2k(N+1)}
 \sum_{r=0}^{n-2k}\frac{\Fb^{(r+2k)}(x)}{r!}h^{r}\,,\;}
 \label{rf:eq:local-cell}
\end{equation}
with the finite reciprocal coefficients $e_{k}$ of
\cref{rf:eq:moments-e-partition}.  Equivalently, on this cell
$G_{N}(x+h)=[M(2^{-N-1}\Dop_{h})]^{-1}P_{x}(h)$, the inverse meaning its
finite action on polynomials of degree at most $n$.
\end{theorem}

\begin{proof}
Write $R=N-n$.  The centred prefix splits into independent parts as
$Y_{N}=X_{n}+2^{-n}Y_{R}$ with $Y_{R}=(1+Z_{R})/2$, $Y_{0}=\tfrac12$, and
$Y_{R}$ is supported in $[2^{-R-1},1-2^{-R-1}]$.  The bound on $h$ gives
$\AbsoluteValue{2^{n}h}\le2^{-R-1}$, so in the box formula for the first
$n$ variables the cutoff pattern cannot change: all indices $\ell<m$
contribute full powers and all $\ell\ge m$ contribute zero (at the two cell
endpoints some arguments can equal zero, which changes nothing since
$n\ge1$).  Thus
\begin{equation}
 G_{N}(x+h)=\frac{2^{-\Tr{n}}}{n!}\sum_{\ell=0}^{m-1}\tm{\ell}\,
 \E\bigl(2m-2\ell-1+2^{n+1}h-Z_{R}\bigr)^{n}.
 \label{rf:eq:local-tail-polynomial}
\end{equation}
Replacing $Z_{R}$ by $Z$ produces a polynomial in $h$ whose $r$th derivative
at $0$ is \cref{rf:eq:splines-derivative-master}; that polynomial is
exactly $P_{x}(h)$.  For any polynomial $p(v)$, symmetry gives
$\E p(v-Z_{R})=M_{R}(\Dop_{v})p(v)$ and $\E p(v-Z)=M(\Dop_{v})p(v)$
(\cref{rf:lem:splines-local-inverse} with $\eta=1$, or directly by Taylor
expansion), and $M_{R}(t)=M(t)/M(2^{-R}t)$.  These identities are
legitimate in the finite-dimensional space of polynomials.  In
\cref{rf:eq:local-tail-polynomial} the variable is
$v=2m-2\ell-1+2^{n+1}h$, so $\Dop_{v}=2^{-n-1}\Dop_{h}$ and the denominator
is $M(2^{-R-n-1}\Dop_{h})=M(2^{-N-1}\Dop_{h})$.  Expanding its reciprocal
by \cref{rf:eq:moments-inverse-M} gives \cref{rf:eq:local-cell}.
\end{proof}

\begin{corollary}[Polynomial dependence on the squared tail scale]
\label{rf:cor:scale-polynomial}
For a fixed dyadic $x=m/2^{n}\in[0,1]$, $n\ge1$, $d=\Floor{n/2}$, and all
integers $N\ge n$,
\begin{equation}
 G_{N}(x)=Q_{x}(4^{-N}),\qquad
 Q_{x}(z)\DefinitionEquals\sum_{k=0}^{d}e_{k}\,4^{-k}\,\Fb^{(2k)}(x)\,z^{k},
 \label{rf:eq:scale-polynomial}
\end{equation}
so $Q_{x}(0)=\Fb(x)$ and $\deg Q_{x}\le d$.
\end{corollary}

This is an exact identity for the whole tail of the sequence
$(G_{N}(x))_{N\ge n}$, not a finite-order asymptotic expansion; its constant
term can therefore be extracted from finitely many depths.  At $x=\tfrac14$,
where $\Fb'=1$, $\Fb''=8$ and the higher derivatives vanish,
\begin{equation}
 G_{N}\!\left(\tfrac14+h\right)=\frac{5}{72}+h+4h^{2}-\frac19\,4^{-N}
 \qquad(N\ge2,\ \AbsoluteValue{h}\le2^{-N-1}),
 \label{rf:eq:quarter-cell}
\end{equation}
which, after the index shift of \cref{rf:sec:extraction-volume}, is the
quarter-point cell that is kernel-verified in the repository
(\cref{rf:sec:register}).

\begin{remark}[The degree drop seen through Prouhet cancellation]
\label{rf:rmk:splines-prouhet}
That the polynomial piece of $G_{N}$ on the cell has degree at most $n$,
although a general piece has degree $N$, can also be seen combinatorially.
The finite product \cref{rf:eq:binary-tm-product} has a zero of order $L$
at $z=1$, so differentiating it at $z=1$ gives Prouhet's cancellation
\begin{equation}
 \sum_{r=0}^{2^{L}-1}\tm{r}\,r^{i}=0\qquad(0\le i<L).
 \label{rf:eq:splines-prouhet}
\end{equation}
With $N=n+L$ and the active indices $h=b2^{L}+r$, $0\le b<m$, $0\le r<2^{L}$,
the signs factor as $\tm{h}=\tm{b}\tm{r}$, and each block of the cell
polynomial is a signed sum of powers $(A+2^{n+1}h'-Br)^{n+L}$ in the cell
variable $h'$; the coefficient of $h'^{i}$ is a polynomial in $r$ of degree
$n+L-i$, which is annihilated by \cref{rf:eq:splines-prouhet} when $i>n$.
This gives the degree bound without the tail; the deconvolution then
identifies the coefficients.
\end{remark}

\subsection{Why the degree bound is attained}
\label{rf:sec:splines-degree}

\begin{lemma}[Signs of the reciprocal coefficients]
\label{rf:lem:splines-e-sign}
For every $k\ge0$, $(-1)^{k}e_{k}>0$.
\end{lemma}

\begin{proof}
Euler's product for the sine \cite[Eq.~4.22.1]{dlmf} gives, near $t=0$,
\begin{equation}
 \frac{1}{M(\ImaginaryUnit t)}
 =\prod_{j=1}^{\infty}\prod_{\ell=1}^{\infty}
 \Bigl(1-\frac{t^{2}}{4^{j}\pi^{2}\ell^{2}}\Bigr)^{-1}.
 \label{rf:eq:splines-euler-product}
\end{equation}
The positive numbers $(4^{j}\pi^{2}\ell^{2})^{-1}$ have finite sum, so the
product converges near zero and its power-series coefficients are limits of
those of finite products.  Every coefficient of a finite product of
geometric series $\sum_{i}(\alpha t^{2})^{i}$ with $\alpha>0$ is
nonnegative, and for each degree $k$ even one factor contributes a strictly
positive coefficient; hence the coefficient of $t^{2k}$ on the left, which
is $(-1)^{k}e_{k}$, is strictly positive.  The infinite product is used
only for the sign; the evaluation formula for $e_{k}$ remains the finite
\cref{rf:eq:moments-e-partition}.
\end{proof}

\begin{proposition}[Exact degree at a reduced interior dyadic]
\label{rf:prop:exact-degree}
If $n\ge1$, $0<m<2^{n}$ and $m$ is odd, then $\deg Q_{m/2^{n}}=\Floor{n/2}$.
\end{proposition}

\begin{proof}
By \cref{rf:eq:setup-global-derivative} and
\cref{rf:eq:setup-integer-values}: if $n=2d$, the top even derivative is
$\Fb^{(2d)}(m/2^{2d})=2^{\Tr{2d}}\Fg(m)=2^{\Tr{2d}}\tm{(m-1)/2}\neq0$; if
$n=2d+1$, it is $\Fb^{(2d)}(m/2^{2d+1})=2^{\Tr{2d}}\Fg(m/2)
=2^{\Tr{2d}-1}\tm{\Floor{m/4}}\neq0$, which also covers $n=1$, $d=0$.  The
coefficient of $z^{d}$ in \cref{rf:eq:scale-polynomial} is then nonzero by
\cref{rf:lem:splines-e-sign}.
\end{proof}

A dyadic argument should therefore be reduced before choosing the number of
levels: an unreduced depth gives a correct formula with a correct upper
bound, but introduces unnecessary interpolation nodes.

% =====================================================================
\section{Quarter-base extraction and its sharpness}
\label{rf:sec:extraction}

\subsection{Extraction from arbitrary distinct depths}
\label{rf:sec:extraction-arbitrary}

\begin{theorem}[Extraction from arbitrary depths]
\label{rf:thm:extraction-arbitrary}
Fix $n\ge1$, $d=\Floor{n/2}$, and distinct integers $N_{0},\ldots,N_{d}\ge n$.
For every $x=m/2^{n}\in[0,1]$,
\begin{equation}
 \Fb(x)=\sum_{i=0}^{d}W_{i}\,G_{N_{i}}(x),\qquad
 W_{i}=\prod_{\substack{0\le\ell\le d\\ \ell\neq i}}
 \frac{4^{-N_{\ell}}}{4^{-N_{\ell}}-4^{-N_{i}}}.
 \label{rf:eq:extraction-arbitrary}
\end{equation}
The weights depend only on the selected depths, not on $m$.
\end{theorem}

\begin{proof}
Let $z_{i}=4^{-N_{i}}$.  For any polynomial $Q$ of degree at most $d$,
$Q(z)=\sum_{i}Q(z_{i})\prod_{\ell\neq i}(z-z_{\ell})/(z_{i}-z_{\ell})$: the
difference of the two sides has degree at most $d$ and vanishes at the $d+1$
distinct points $z_{i}$, so it is zero.  Evaluate at $z=0$ with $Q=Q_{x}$
from \cref{rf:cor:scale-polynomial}.
\end{proof}

\subsection{Consecutive depths and the Gaussian-binomial row}
\label{rf:sec:extraction-row}

For $0<q<1$ and $0\le j\le d$ define
\begin{equation}
 w_{d,j}(q)\DefinitionEquals
 \prod_{\substack{0\le\ell\le d\\ \ell\neq j}}\frac{-q^{\ell}}{q^{j}-q^{\ell}},
 \label{rf:eq:extraction-weights-product}
\end{equation}
and $w_{d,j}\DefinitionEquals w_{d,j}(\tfrac14)$.

\begin{lemma}[Constant extraction on a geometric grid]
\label{rf:lem:extraction-weights}
For every polynomial $P$ of degree at most $d$ and every $z\neq0$,
\begin{equation}
 P(0)=\sum_{j=0}^{d}w_{d,j}(q)\,P(zq^{j}),
 \label{rf:eq:extraction-filter}
\end{equation}
equivalently $\sum_{j=0}^{d}w_{d,j}(q)q^{jr}=[r=0]$ for $0\le r\le d$.  In
closed form,
\begin{equation}
 \boxed{\;
 w_{d,j}(q)=\frac{(-1)^{d-j}q^{\Tr{d-j}}}{\qp{q}{q}{j}\qp{q}{q}{d-j}}
 =\frac{(-1)^{d-j}q^{\Tr{d-j}}}{\qp{q}{q}{d}}\qb{d}{j}{q}\,,\;}
 \label{rf:eq:extraction-weights-q}
\end{equation}
and the generating polynomial of the row is
\begin{equation}
 \sum_{j=0}^{d}w_{d,j}(q)\,v^{j}=\prod_{r=1}^{d}\frac{v-q^{r}}{1-q^{r}}.
 \label{rf:eq:extraction-row-polynomial}
\end{equation}
At $q=\tfrac14$, $\qp{q}{q}{r}=4^{-\Tr{r}}\prod_{s=1}^{r}(4^{s}-1)$, which
turns \cref{rf:eq:extraction-weights-q} into the integer-product form
\cref{rf:eq:main-weights-integer}.
\end{lemma}

\begin{proof}
The Lagrange basis polynomial for the node $zq^{j}$ among the nodes
$zq^{0},\ldots,zq^{d}$, evaluated at $0$, is
$\prod_{\ell\neq j}(0-zq^{\ell})/(zq^{j}-zq^{\ell})$, in which $z$ cancels;
that is \cref{rf:eq:extraction-weights-product}, and Lagrange interpolation
(the uniqueness argument of the previous proof) gives
\cref{rf:eq:extraction-filter}; the monomials $P=z^{r}$ give the moment
identities.  Splitting the product into $\ell<j$, where the factor is
$1/(1-q^{j-\ell})$, and $\ell>j$, where it is $-q^{\ell-j}/(1-q^{\ell-j})$,
gives \cref{rf:eq:extraction-weights-q}, the exponent being
$\sum_{i=1}^{d-j}i=\Tr{d-j}$.  The generating polynomial has degree $d$,
vanishes at $v=q^{r}$ for $1\le r\le d$ by the moment identities, and takes
the value $1$ at $v=1$ (the case $r=0$), which determines it as the right
side of \cref{rf:eq:extraction-row-polynomial}.  The closed forms are also
consequences of the finite $q$-binomial theorem \cite[\S17.2]{dlmf}, but no
such theorem is needed.
\end{proof}

\begin{theorem}[Extraction at consecutive depths]
\label{rf:thm:extraction-consecutive}
For $n\ge1$, $0\le m\le2^{n}$, $d=\Floor{n/2}$ and every integer $R\ge0$,
\begin{equation}
 \boxed{\;
 \Fb\!\left(\frac{m}{2^{n}}\right)
 =\sum_{j=0}^{d}w_{d,j}\,G_{n+R+j}\!\left(\frac{m}{2^{n}}\right),\;}
 \label{rf:eq:extraction-consecutive}
\end{equation}
and substituting \cref{rf:eq:splines-G-integer} gives
\cref{rf:eq:main-quarter-shift}, hence \cref{rf:eq:main-quarter} at $R=0$.
In fully explicit $q$-notation, with $Q=\tfrac14$,
\begin{equation}
 \boxed{\begin{aligned}
 \Fb\!\left(\frac{m}{2^{n}}\right)={}&\frac{1}{\qp{Q}{Q}{d}}\sum_{j=0}^{d}
 \frac{(-1)^{d-j}Q^{\Tr{d-j}}}{2^{\Tr{n+j}}\,(n+j)!}\qb{d}{j}{Q}\\[-0.2em]
 &\times\sum_{h=0}^{2^{j}m-1}\tm{h}\,\bigl(2^{j+1}m-2h-1\bigr)^{n+j}.
 \end{aligned}}
 \label{rf:eq:extraction-closed}
\end{equation}
No moments, Bernoulli numbers, derivative values, recursively defined
coefficients, limits, or coefficient extractors occur in this formula.
\end{theorem}

\begin{proof}
Apply \cref{rf:lem:extraction-weights} to the polynomial $Q_{x}$ of
\cref{rf:cor:scale-polynomial} with $q=\tfrac14$ and $z=4^{-(n+R)}$, whose
nodes are the depths $n+R+j$; equivalently, this is
\cref{rf:thm:extraction-arbitrary} with $N_{j}=n+R+j$, the common factor
$4^{-(n+R)}$ cancelling from the weights.  Then replace each
$G_{n+R+j}(m/2^{n})$ by the integer sum \cref{rf:eq:splines-G-integer} with
$A=2^{R+j}m$, and at $R=0$ insert \cref{rf:eq:extraction-weights-q}.  The
case $d=0$ uses the single weight one.
\end{proof}

The row \cref{rf:eq:main-quarter} uses $d+1$ splines of orders
$n,\ldots,n+d$; its largest power is $n+d$, and before simplification it
contains $\sum_{j=0}^{d}2^{j}m=m(2^{d+1}-1)$ signed integer powers.  It is
attractive when the absence of every auxiliary coefficient sequence is more
important than minimizing arithmetic; for large numerators the block
formula is usually the better organization.  Consecutive depths need not
start at $n$ (the parameter $R$), and arbitrary distinct depths $\ge n$ are
covered by \cref{rf:thm:extraction-arbitrary}, which permits the reuse of
already available spline values.

\subsection{A precise optimality statement}
\label{rf:sec:extraction-minimal}

The number $d+1$ is not a lower bound on all possible algorithms --- the
binary formulas are not interpolation rules --- but it is sharp within the
natural class of grid-wide linear rules.

\begin{theorem}[Minimal number of samples for a grid-wide linear rule]
\label{rf:thm:extraction-minimal}
Fix $n\ge1$ and $d=\Floor{n/2}$.  Suppose distinct depths
$N_{1},\ldots,N_{L}\ge n$ and real numbers $c_{1},\ldots,c_{L}$, independent
of $m$, satisfy $\Fb(m/2^{n})=\sum_{i=1}^{L}c_{i}G_{N_{i}}(m/2^{n})$ for
every $m=0,\ldots,2^{n}$.  Then $L\ge d+1$.
\end{theorem}

\begin{proof}
The grid contains $x_{0}=1$ and $x_{r}=2^{-2r}$ for $1\le r\le d$.  The
polynomial $Q_{x_{0}}$ is the constant $1$, and by
\cref{rf:prop:exact-degree} $Q_{x_{r}}$ has degree exactly $r$.  These
$d+1$ polynomials therefore form a basis of the polynomials of degree at
most $d$: their leading terms make the coefficient matrix triangular with
nonzero diagonal.  The supposed rule recovers $Q(0)$ from the values
$Q(z_{i})$, $z_{i}=4^{-N_{i}}>0$, for each basis element, hence by
linearity for every polynomial $Q$ of degree at most $d$.  If $L\le d$,
take $Q(z)=\prod_{i=1}^{L}(z-z_{i})$: all sampled values vanish but
$Q(0)\neq0$, a contradiction; $L=0$ is included by the empty product.
\end{proof}

The theorem does not exclude weights depending on the point, nonlinear
formulas, other kinds of samples, or a direct finite coefficient
calculation.

\subsection{The extraction weights are well conditioned}
\label{rf:sec:extraction-stability}

\begin{proposition}[Absolute row sum]
\label{rf:prop:extraction-norm}
For $0<q<1$,
\begin{equation}
 \sum_{j=0}^{d}\AbsoluteValue{w_{d,j}(q)}
 =\frac{\qp{-q}{q}{d}}{\qp{q}{q}{d}}
 =\prod_{r=1}^{d}\frac{1+q^{r}}{1-q^{r}}
 \le\prod_{r=1}^{\infty}\frac{1+q^{r}}{1-q^{r}}<\infty.
 \label{rf:eq:extraction-norm}
\end{equation}
For $q=\tfrac14$ the absolute sum is strictly less than $2$, uniformly in
$d$; the limiting values are $1.96926\ldots$ at $q=\tfrac14$ and
$8.25599\ldots$ at $q=\tfrac12$.  No nonnegative extraction row exists when
$d\ge1$.
\end{proposition}

\begin{proof}
The sign of $w_{d,j}(q)$ is $(-1)^{d-j}$, so the absolute sum is
$(-1)^{d}\sum_{j}w_{d,j}(q)(-1)^{j}$, which by
\cref{rf:eq:extraction-row-polynomial} at $v=-1$ equals
$\prod_{r}(1+q^{r})/(1-q^{r})$; this is the first equality, and
$\prod_{r\le d}(1+q^{r})=\qp{-q}{q}{d}$.  The infinite product is finite
because $\log((1+u)/(1-u))\le2u/(1-u)$ for $0\le u<1$ and $\sum_{r}q^{r}$
converges.  For $q=\tfrac14$ the first factor is $\tfrac53$ and
$\sum_{r\ge2}\log\frac{1+4^{-r}}{1-4^{-r}}\le\frac{32}{15}\sum_{r\ge2}4^{-r}=\frac{8}{45}$,
while $\log\frac65\ge\frac15-\frac1{50}=\frac{9}{50}>\frac{8}{45}$; hence the
product is below $\tfrac53\cdot\tfrac65=2$ for every $d$ (for $d=0$ the sum
is $1$).  Finally, positive nodes with nonnegative weights summing to one
cannot have a vanishing first moment, so a nonnegative row is impossible
when $d\ge1$.
\end{proof}

If each exactly computed input $G_{N_{j}}(x)$ carries an absolute error at
most $\delta$, combining them with the exact weights adds error below
$2\delta$.  This bounds the amplification by the outer filter only; the
alternating power sum inside an individual spline can cancel severely, and
exact integer or rational arithmetic is the appropriate default there.

\subsection{Dictionary to the canonical extraction volume}
\label{rf:sec:extraction-volume}

The canonical volume of this subgroup \cite{proveit-due} works with the
density splines $\up_{n}$, the inverse Fourier transforms of the prefix
products with factors $k=0,\ldots,n$ of
$\prod_{k\ge0}\SincRad(\pi\xi/2^{k})$, and with the reciprocal-product
coefficients $\amvol{r}$ defined by $1/\FourierTwoPi{\up}(\xi)=\sum_{r}\amvol{r}(2\pi\xi)^{2r}$.
In the notation of this document,
\begin{equation}
 \up_{n}=u_{n+1},\qquad
 \amvol{r}=(-1)^{r}e_{r},\qquad
 G_{N}(x)=\up_{N-1}\text{'s distribution function}=\Fb\text{-side of }u_{N},
 \label{rf:eq:extraction-dictionary}
\end{equation}
where $u_{N}$ is the density of $Z_{N}$ (\cref{rf:sec:rvachev-density});
indeed $\FourierTwoPi{\up}(\xi)=M(2\pi\ImaginaryUnit\xi)$, so
$1/\FourierTwoPi{\up}(\xi)=\sum_{k}e_{k}(2\pi\ImaginaryUnit\xi)^{2k}
=\sum_{k}(-1)^{k}e_{k}(2\pi\xi)^{2k}$, and the volume's
$\amvol{1}=\tfrac1{18}$, $\amvol{2}=\tfrac{31}{16200}$, $\amvol{3}=\tfrac{2347}{42865200}$
are the absolute values of \cref{rf:eq:moments-e-values}.  The volume's main
theorem, $\up_{n}(x)=\up(x)+\sum_{r=1}^{d}(-1)^{r}\amvol{r}4^{-r}\up^{(2r)}(x)4^{-rn}$
for $n\ge s(x)$, is \cref{rf:eq:rvachev-density-law} below read through this
dictionary: $u_{N}(x)=\sum_{k}e_{k}4^{-Nk}\up^{(2k)}(x)$ with $N=n+1$.  Its
extraction row, with the weights
$(-1)^{d-j}Q^{\binom{d-j+1}{2}}\qb{d}{j}{Q}/\qp{Q}{Q}{d}$, is the row
\cref{rf:eq:extraction-weights-q} of the present document, and its depth
$s(x)$ of a dyadic $x\in[-1,1]$ is the reduced denominator exponent used
here.  The volume also proves, from Rvachev's functional-differential
equation, that the defect at the last level before the onset is
$-\tm{k}2^{-\binom{s}{2}}B_{s}/s!$, and records which pieces are
kernel-verified; \cref{rf:sec:register} collects the crosswalks relevant
to the formulas of this document.

% =====================================================================
\section{The half-base formula and the general shift}
\label{rf:sec:halfbase}

The quarter-base formula is not obtained by replacing $\tfrac12$ by
$\tfrac14$ in an existing identity: centring removes all odd cutoff powers
and halves the interpolation degree.  A shift parameter makes this precise.
For $0\le\alpha\le1$ set
\begin{equation}
 G_{N,\alpha}(x)\DefinitionEquals\Prob(X_{N}+\alpha2^{-N}\le x),
 \label{rf:eq:halfbase-shifted}
\end{equation}
so $G_{N,1/2}=G_{N}$ and $G_{N,0}=H_{N}$.  For $x=m/2^{n}$, $N\ge n\ge1$ and
$A=2^{N-n}m$, the box formula gives
\begin{equation}
 G_{N,\alpha}\!\left(\frac{m}{2^{n}}\right)
 =\frac{2^{-\binom{N}{2}}}{N!}\sum_{h=0}^{A-1}\tm{h}\,(A-h-\alpha)^{N};
 \label{rf:eq:halfbase-shifted-finite}
\end{equation}
at $\alpha\in\SetOf{0,1}$ the extra boundary terms have vanishing $N$th
power, so the same expression covers both endpoints.  For arbitrary complex
$\alpha$ define $G_{N,\alpha}(m/2^{n})$ by the polynomial on the right; it
loses its probability interpretation but remains a finite algebraic
object.

\begin{proposition}[General shift and the half-base row]
\label{rf:prop:halfbase}
For $n\ge1$, $0\le m\le2^{n}$ and every $\alpha\in\ComplexNumbers$,
\begin{equation}
 \Fb\!\left(\frac{m}{2^{n}}\right)
 =\sum_{j=0}^{n}w_{n,j}\!\left(\tfrac12\right)G_{n+j,\alpha}\!\left(\frac{m}{2^{n}}\right),
 \label{rf:eq:halfbase-extract}
\end{equation}
equivalently
\begin{equation}
 \boxed{\begin{aligned}
 \Fb\!\left(\frac{m}{2^{n}}\right)={}&\frac{1}{2^{n^{2}}\qp{\frac12}{\frac12}{n}}
 \sum_{j=0}^{n}\frac{\qb{n}{j}{1/2}}{2^{j(j-1)}\,(n+j)!}\\[-0.1em]
 &\times\sum_{h=0}^{2^{j}m-1}\tm{h}\,\bigl(h-2^{j}m+\alpha\bigr)^{n+j}.
 \end{aligned}}
 \label{rf:eq:halfbase-reshetnikov}
\end{equation}
\end{proposition}

\begin{proof}
First let $0\le\alpha\le1$ and $\eta=2^{-N}$.  The head--tail identity
reads $\Fb(t)=\E G_{N,\alpha}(t-\eta(X'-\alpha))$, and at $x=A\eta$ the
whole averaging interval $[x-\eta(1-\alpha),x+\eta\alpha]$ lies in one
polynomial cell of $G_{N,\alpha}$, whose endpoints are consecutive possible
knots.  Apply \cref{rf:lem:splines-local-inverse} with $W=X'-\alpha$
(supported in $[-1,1]$, not symmetric) and the degree $N$: since the jet of
$\Fb$ at $x$ terminates at order $n$,
\[
 G_{N,\alpha}(x)=\sum_{r=0}^{n}\gamma_{r}(\alpha)\,\eta^{r}\Fb^{(r)}(x),\qquad
 \gamma_{r}(\alpha)=\coeff{z}{r}\frac{1}{\MomentGeneratingFunction{W}(-z)},
\]
with $\gamma_{0}=1$ and coefficients independent of $N$.  Thus
$G_{N,\alpha}(x)$ is an exact polynomial of degree at most $n$ in $2^{-N}$
with constant term $\Fb(x)$, and \cref{rf:lem:extraction-weights} at
$q=\tfrac12$ with $n+1$ consecutive depths gives
\cref{rf:eq:halfbase-extract}.  When \cref{rf:eq:halfbase-shifted-finite}
is substituted, both sides are polynomials in $\alpha$; equality on $[0,1]$
implies equality for all complex $\alpha$.  Finally insert the weights
\cref{rf:eq:extraction-weights-q} at $q=\tfrac12$ and write
$(A-h-\alpha)^{n+j}=(-1)^{n+j}(h-A+\alpha)^{n+j}$, which cancels the sign
$(-1)^{n-j}$ of the weight; the powers of two combine as
$\Tr{n-j}+\binom{n+j}{2}=n^{2}+j(j-1)$, giving
\cref{rf:eq:halfbase-reshetnikov}.
\end{proof}

Formula \cref{rf:eq:halfbase-reshetnikov} is the unit-interval form of the
expression posed by Reshetnikov in 2019 \cite{reshetnikov2019} and
subsequently developed in the repository \cite{proveit-primary}, which also
treats the signed extension at unrestricted numerators; the fold
\cref{rf:eq:rvachev-global-fold} evaluates it everywhere without reproving
that identity.  At $\alpha=\tfrac12$ the variable $W=Z'/2$ is symmetric, the
odd coefficients $\gamma_{r}$ vanish, and the degree drops from at most $n$
in $2^{-N}$ to at most $\Floor{n/2}$ in $4^{-N}$: this is precisely the
passage from \cref{rf:eq:halfbase-extract} to
\cref{rf:eq:extraction-consecutive}, the half-base row using $n+1$ splines
of orders up to $2n$ where the quarter-base row uses $d+1$ of orders up to
$n+d$.

% =====================================================================
\section{An integer determinant, a common denominator, and certified rounding}
\label{rf:sec:determinant}

\subsection{The bordered determinant}
\label{rf:sec:determinant-formula}

The composition and partition formulas expose individual coefficients.  A
bordered integer matrix packages the entire dyadic value into a single
ratio of integers; its Leibniz expansion is a finite sum of products, so no
matrix inversion or recurrence is part of the definition.

Put $d=\Floor{n/2}$ and define the integer row entries
\begin{equation}
 t_{j}\DefinitionEquals\binom{n}{2j}\,\Spow{n-2j}{m}
 =\binom{n}{2j}\sum_{h=0}^{m-1}\tm{h}\,(2m-2h-1)^{n-2j}\qquad(0\le j\le d),
 \label{rf:eq:determinant-row}
\end{equation}
with $\Spow{0}{m}=\sum_{h<m}\tm{h}$.  Let $L_{d}$ be the $d\times d$ lower
triangular integer matrix with indices $1\le k,j\le d$ and entries
\begin{equation}
 (L_{d})_{kj}=
 \begin{cases}
 (2k+1)(4^{k}-1),& j=k,\\
 -\binom{2k+1}{2j},& j<k,\\
 0,& j>k,
 \end{cases}
 \label{rf:eq:determinant-L}
\end{equation}
and form the bordered $(d+1)\times(d+1)$ matrix
\begin{equation}
 B_{n,m}\DefinitionEquals
 \begin{pmatrix}
 L_{d} & \mathbf 1\\
 -t_{1}\ \cdots\ -t_{d} & t_{0}
 \end{pmatrix},
 \label{rf:eq:determinant-B}
\end{equation}
where $\mathbf 1$ is the column of $d$ ones; for $d=0$ this is the
$1\times1$ matrix $(t_{0})$.  Every entry is an explicit integer.

\begin{theorem}[Integer determinant formula]
\label{rf:thm:determinant}
For $n\ge0$ and $0\le m\le2^{n}$,
\begin{equation}
 \boxed{\;
 \Fb\!\left(\frac{m}{2^{n}}\right)
 =\frac{\det B_{n,m}}{D_{n}},\qquad
 D_{n}\DefinitionEquals2^{\Tr{n}}\,n!\prod_{k=1}^{d}(2k+1)(4^{k}-1)
 =2^{\Tr{n}}\,n!\,(2d+1)!!\prod_{k=1}^{d}(4^{k}-1).\;}
 \label{rf:eq:determinant-value}
\end{equation}
In particular the reduced denominator of every value on the grid of spacing
$2^{-n}$, and of every Rvachev value $\up(m/2^{n})$, divides $D_{n}$.
\end{theorem}

\begin{proof}
Multiplying \cref{rf:eq:moments-recurrence} by $(2k+1)!$ and separating the
top term gives, in terms of the moments $c_{j}=(2j)!a_{j}$,
\[
 (2k+1)(4^{k}-1)\,c_{k}=\sum_{j=0}^{k-1}\binom{2k+1}{2j}c_{j}\qquad(k\ge1),
\]
and since $c_{0}=1$ this is the finite system
$L_{d}(c_{1},\ldots,c_{d})^{\mathsf T}=\mathbf 1$.  The diagonal entries are
nonzero and $\det L_{d}=\prod_{k=1}^{d}(2k+1)(4^{k}-1)$.  For a block
matrix with invertible upper-left block, $\det\begin{pmatrix}L&\mathbf 1\\-t^{\mathsf T}&t_{0}\end{pmatrix}
=\det L\cdot\bigl(t_{0}+t^{\mathsf T}L^{-1}\mathbf 1\bigr)$ (the Schur
complement; equivalently, add $t^{\mathsf T}L^{-1}$ times the first $d$ rows
to the last row, which clears its first $d$ entries).  With
$L_{d}^{-1}\mathbf 1=c=(c_{1},\ldots,c_{d})^{\mathsf T}$ this gives
\[
 \det B_{n,m}=\det L_{d}\Bigl(t_{0}+\sum_{j=1}^{d}t_{j}c_{j}\Bigr).
\]
By \cref{rf:eq:master-classical}, $\Fb(m/2^{n})=2^{-\Tr{n}}(n!)^{-1}\sum_{j=0}^{d}t_{j}c_{j}$,
which is \cref{rf:eq:determinant-value}; the case $d=0$ is the same formula
directly.  The moment relation was used only to identify the matrix; the
final expression asks for no moment.  Integrality of the entries gives the
denominator statement for $\Fb$, and the fold \cref{rf:eq:setup-fold} for
$\up$.
\end{proof}

In the literal finite-sum reading,
\begin{equation}
 \det B_{n,m}=\sum_{\sigma\in\mathfrak S_{d+1}}
 (-1)^{\#\SetOf{(i,j):\,i<j,\ \sigma(i)>\sigma(j)}}
 \prod_{i=1}^{d+1}(B_{n,m})_{i,\sigma(i)},
 \label{rf:eq:determinant-leibniz}
\end{equation}
so even the determinant notation can be removed; many terms vanish because
$L_{d}$ is triangular, and an implementation should use exact elimination
rather than this factorial-size expansion.  For $n=2$, $m=1$:
$B_{2,1}=\begin{pmatrix}9&1\\-1&1\end{pmatrix}$, $\det B_{2,1}=10$,
$D_{2}=2^{3}\cdot2!\cdot9=144$, and $\Fb(\tfrac14)=\tfrac{10}{144}=\tfrac5{72}$.

\begin{remark}[An equivalent $(d+2)$-bordering]
\label{rf:rmk:determinant-alt}
One may also border the full $(d+1)\times(d+1)$ triangular matrix $T$ with
$T_{00}=1$, $T_{ii}=(2i+1)(4^{i}-1)$, $T_{ij}=-\binom{2i+1}{2j}$ for $j<i$,
by the unit vector $e_{0}$ and the row $v=(t_{0},\ldots,t_{d})$:
$\Fb(m/2^{n})=-\det\begin{pmatrix}T&e_{0}\\ v^{\mathsf T}&0\end{pmatrix}/D_{n}$,
by the same Schur-complement identity
$\det=-\det T\cdot v^{\mathsf T}T^{-1}e_{0}=-\det T\cdot v^{\mathsf T}c$.
The $(d+1)$-form above absorbs the row $c_{0}=1$ into the border.
\end{remark}

\subsection{A second denominator bound, and the prefix error}
\label{rf:sec:determinant-denominator}

$D_{n}$ is a convenient common denominator, not a claim of minimality, and
it does not replace the sharper arithmetic of \cite{arias2018}.  A different
elementary bound follows directly from the positive composition formula.

\begin{proposition}[Denominator bound from compositions]
\label{rf:prop:determinant-denominator}
For $n\ge0$ and $d=\Floor{n/2}$, the number
$D'_{n}\DefinitionEquals2^{\Tr{n}}(n+d)!\prod_{j=1}^{d}(4^{j}-1)$ satisfies
$D'_{n}\Fb(m/2^{n})\in\IntegerNumbers$ for $0\le m\le2^{n}$.
\end{proposition}

\begin{proof}
Use \cref{rf:eq:composition-value}.  A summand with index $k\le d$ and
composition $k=r_{1}+\cdots+r_{\ell}$ has denominator dividing
$2^{\Tr{n}}(n-2k)!\prod_{i}(2r_{i}+1)!\prod_{i}(4^{R_{i}}-1)$.  The suffix
sums $R_{i}$ are distinct integers in $[1,k]$, so the last product divides
$\prod_{j\le d}(4^{j}-1)$.  The factorial arguments add up to
$(n-2k)+\sum_{i}(2r_{i}+1)=n+\ell\le n+d$, and by integrality of
multinomial coefficients their product divides $(n+\ell)!$, which divides
$(n+d)!$; the term $k=0$ has just $n!$.  Every summand, hence their
integer-signed sum, has denominator dividing $D'_{n}$.
\end{proof}

Neither bound dominates the other in general: $D_{4}=16588800$ while
$D'_{4}=2^{10}\cdot720\cdot45=33177600$, but $D'_{n}$ has the smaller
factorial for large $n$ relative to the double factorial.  Either may be
used below.

\begin{lemma}[Uniform prefix error]
\label{rf:lem:determinant-prefix-error}
For every $N\ge1$ and every real $x$,
$\AbsoluteValue{G_{N}(x)-\Fb(x)}\le2^{-N}$.
\end{lemma}

\begin{proof}
Couple the variables by $X=Y_{N}+2^{-N-1}Z'$ with $Z'$ an independent copy
of $Z$ and $\AbsoluteValue{Z'}\le1$.  With $\delta=2^{-N-1}$ this gives
$\Fb(x-\delta)\le G_{N}(x)\le\Fb(x+\delta)$, and the $2$-Lipschitz bound of
\cref{rf:lem:setup-normalization} proves the claim.
\end{proof}

\begin{corollary}[Certified one-prefix rounding]
\label{rf:cor:rounding}
For $n\ge1$ let $D$ be any of the common denominators $D_{n}$, $D'_{n}$,
and choose the explicit integer $N=\max\SetOf{n,\ 2+\Floor{\log_{2}D}}$, so
that $2^{N}>2D$.  Then
\begin{equation}
 \boxed{\;
 \Fb\!\left(\frac{m}{2^{n}}\right)
 =\frac{1}{D}\Floor{\frac{D\,\Spow{N}{2^{N-n}m}}{2^{\Tr{N}}N!}+\frac12}\,.\;}
 \label{rf:eq:rounding}
\end{equation}
\end{corollary}

\begin{proof}
By \cref{rf:lem:determinant-prefix-error},
$\AbsoluteValue{D\,G_{N}(x)-D\,\Fb(x)}\le D2^{-N}<\tfrac12$, and $D\Fb(x)$
is an integer; nearest-integer rounding is therefore exact, with no tie.
Substitute \cref{rf:eq:splines-G-integer}.
\end{proof}

The logarithm in the choice of $N$ can be replaced by the binary digit
length of the explicitly specified integer $D$.  The formula has only one
finite prefix, but that prefix is enormous; it is a mathematical
certification device, not a recommended implementation.

% =====================================================================
\section{Rvachev values, the signed extension, derivatives and primitives}
\label{rf:sec:rvachev}

\subsection{Every dyadic bump value}
\label{rf:sec:rvachev-fold}

For an integer $a$ and $n\ge0$, the support and the fold
\cref{rf:eq:setup-fold} give
\begin{equation}
 \boxed{\;
 \up\!\left(\frac{a}{2^{n}}\right)=
 \begin{cases}
 \Fb\!\left(\dfrac{2^{n}-\AbsoluteValue{a}}{2^{n}}\right), & \AbsoluteValue{a}\le2^{n},\\[2mm]
 0, & \AbsoluteValue{a}>2^{n},
 \end{cases}\;}
 \label{rf:eq:rvachev-fold}
\end{equation}
so every formula for $\Fb$ gives a closed rational formula for $\up$ with
no new coefficients; in particular $\up(0)=1$, $\up(\pm1)=0$,
$\up(\pm\tfrac12)=\tfrac12$, and $\up(\tfrac34)=\Fb(\tfrac14)=\tfrac5{72}$.

\subsection{Direct extraction from finite densities}
\label{rf:sec:rvachev-density}

When the available data are finite sinc products rather than distribution
functions, a direct formula is useful.  Let $u_{N}$ be the density of
$Z_{N}=\sum_{j=1}^{N}2^{-j}V_{j}$; under the $2\pi$ convention
\begin{equation}
 \FourierTwoPi{u_{N}}(\xi)=\prod_{k=0}^{N-1}\SincRad\!\left(\frac{\pi\xi}{2^{k}}\right),
 \qquad\SincRad(z)=\frac{\sin z}{z},
 \label{rf:eq:rvachev-sinc}
\end{equation}
so a product ending with the factor $k=n$ is $u_{n+1}$ in this notation; this
one-factor offset matters in exact cutoff formulas
(\cref{rf:sec:extraction-volume}).  Since
$Z_{N}=\sum_{j=1}^{N}2^{1-j}U_{j}-(1-2^{-N})$, differentiating the
weighted-uniform distribution function of \cref{rf:lem:master-box} once
gives, for $N\ge2$ (so that the density is continuous),
\begin{equation}
 u_{N}(x)=\frac{2^{N(N-1)/2}}{(N-1)!}\sum_{h=0}^{2^{N}-1}\tm{h}
 \PositivePart{x+1-2^{-N}-2^{1-N}h}^{N-1},
 \label{rf:eq:rvachev-density-spline}
\end{equation}
whose knots are the odd multiples of $2^{-N}$; for $N=1$ the same display
with the convention $s_{+}^{0}=[s>0]$ is the indicator density of
$[-\tfrac12,\tfrac12]$.  For $\AbsoluteValue{a}<2^{n}$ and $N>n\ge0$ put
$A=2^{N}+a2^{N-n}$, a positive even integer; the terms with $2h<A-1$ are
exactly the active ones, and
\begin{equation}
 \boxed{\;
 u_{N}\!\left(\frac{a}{2^{n}}\right)
 =\frac{1}{(N-1)!\,2^{N(N-1)/2}}\sum_{h=0}^{A/2-1}\tm{h}\,(A-1-2h)^{N-1}.\;}
 \label{rf:eq:rvachev-density-integer}
\end{equation}
The two finite objects are related by an exact fold: the independent-copy
identity $Z_{N}\EqualInLaw(V+Z'_{N-1})/2$ gives
$u_{N}(x)=\Prob(2x-1\le Z'_{N-1}\le2x+1)$, and for $x\ge0$ the upper cutoff
may be discarded, so by symmetry
\begin{equation}
 u_{N}(x)=G_{N-1}\bigl(1-\AbsoluteValue{x}\bigr)\qquad(x\in\RealNumbers,\ N\ge2).
 \label{rf:eq:rvachev-finite-fold}
\end{equation}

\begin{theorem}[Direct quarter-base extraction for the bump]
\label{rf:thm:up-quarter}
Let $n\ge0$, $\AbsoluteValue{a}<2^{n}$, $d=\Floor{n/2}$, and $x=a/2^{n}$.
Then for every $N>n$
\begin{equation}
 u_{N}(x)=\sum_{k=0}^{d}e_{k}\,4^{-Nk}\,\up^{(2k)}(x),
 \label{rf:eq:rvachev-density-law}
\end{equation}
and for every $R\ge0$
\begin{equation}
 \boxed{\;
 \up\!\left(\frac{a}{2^{n}}\right)
 =\sum_{j=0}^{d}w_{d,j}\,u_{n+1+R+j}\!\left(\frac{a}{2^{n}}\right),\;}
 \label{rf:eq:up-quarter}
\end{equation}
with each $u_{N}$ the finite integer sum \cref{rf:eq:rvachev-density-integer}.
For $\AbsoluteValue{a}\ge2^{n}$ the bump value is zero.
\end{theorem}

\begin{proof}
With $\eta=2^{-N}$, $Z\EqualInLaw Z_{N}+\eta Z'$ gives $\up(t)=\E u_{N}(t-\eta Z')$.
Since $N>n$, the number $x/\eta=a2^{N-n}$ is even while the knots of $u_{N}$
are odd multiples of $\eta$, so $[x-\eta,x+\eta]$ is one polynomial cell, on
which $u_{N}$ has degree $N-1$.  Apply \cref{rf:lem:splines-local-inverse}
with this degree, $W=Z'$ and $\eta$: the jet of $\up$ at $x$ terminates at
order $n$ (\cref{rf:lem:setup-dilation}), only even orders survive by
symmetry, and the reciprocal operator $[M(-\eta\Dop)]^{-1}=\sum_{k}e_{k}\eta^{2k}\Dop^{2k}$
evaluated at $s=0$ gives \cref{rf:eq:rvachev-density-law}.  Constant
extraction on the geometric grid $4^{-N}$, $N=n+1+R+j$, with $q=\tfrac14$
yields \cref{rf:eq:up-quarter}.  The case $N=1$, $n=0$, $a=0$ is the
constant density on its interior and follows by the same degree-zero
argument; outside $[-1,1]$ and at its endpoints the bump vanishes.
Alternatively, apply \cref{rf:thm:extraction-consecutive} at
$(2^{n}-\AbsoluteValue{a})/2^{n}$ with depths $n+R+j$ and fold back by
\cref{rf:eq:rvachev-finite-fold}, which explains the shift by one in the
levels.
\end{proof}

At $x=\tfrac14$: $u_{3}(\tfrac14)=\tfrac{15}{16}$, $u_{4}(\tfrac14)=\tfrac{179}{192}$,
and $\up(\tfrac14)=\bigl(4u_{4}-u_{3}\bigr)/3=\tfrac{67}{72}$; indeed
$\up''(\tfrac14)=-8$, so the whole tail satisfies
$u_{N}(\tfrac14)=\tfrac{67}{72}+\tfrac49\,4^{-N}$ for $N\ge3$.  Two suitable
consecutive values determine the limit with zero remainder, unlike
numerical extrapolation from an asymptotic expansion.

\subsection{The signed global extension at every dyadic argument}
\label{rf:sec:rvachev-global}

For $a\ge0$ define
\begin{equation}
 b\DefinitionEquals\Floor{\frac{a}{2^{n+1}}},\qquad
 r\DefinitionEquals a-2^{n+1}b,\qquad
 a^{\ast}\DefinitionEquals\min\SetOf{r,\,2^{n+1}-r}.
 \label{rf:eq:rvachev-global-data}
\end{equation}
Then $0\le a^{\ast}\le2^{n}$ and, by \cref{rf:eq:setup-global-block} and the
fold of the bump around its centre,
\begin{equation}
 \boxed{\;
 \Fg\!\left(\frac{a}{2^{n}}\right)=\tm{b}\,\Fb\!\left(\frac{a^{\ast}}{2^{n}}\right),\;}
 \qquad\Fg(x)=0\ \ (x<0).
 \label{rf:eq:rvachev-global-fold}
\end{equation}
For instance $\Fb(3)=1$ whereas $\Fg(3)=-1$.  The block formula
\cref{rf:eq:binary-block} performs this reduction implicitly through the
high-bit signs $\eta_{j}$ and needs no separate fold.  No formula requires
$a/2^{n}$ to be reduced: directly from the proved identities any of the
closed expressions satisfies $\mathcal F_{n}(a)=\mathcal F_{n+1}(2a)$, and
reducing merely shortens the sums; the reflection
$\Fb(m/2^{n})+\Fb((2^{n}-m)/2^{n})=1$ is a further useful reduction.

\subsection{Closed forms for every dyadic derivative}
\label{rf:sec:rvachev-derivatives}

For $a,n\ge0$, \cref{rf:eq:setup-global-derivative} and
\cref{rf:eq:setup-integer-values} give
\begin{equation}
 \Fg^{(r)}\!\left(\frac{a}{2^{n}}\right)=
 \begin{cases}
 2^{\Tr{r}}\,\Fg\!\left(a/2^{n-r}\right), & 0\le r\le n,\\
 0, & r\ge n+1,
 \end{cases}
 \label{rf:eq:rvachev-derivatives-global}
\end{equation}
since for $r\ge n+1$ the argument $2^{r}a/2^{n}$ is an even integer.  On the
interior of $[0,1]$ the derivatives of $\Fb$ agree with those of $\Fg$, and
at $0$ and $1$ all positive-order derivatives of $\Fb$ vanish.  For the bump
use the signed extension without an absolute-value fold:
\begin{equation}
 \up^{(r)}(x)=2^{\Tr{r}}\,\Fg\bigl(2^{r}(x+1)\bigr)\qquad(-1\le x\le1),
 \label{rf:eq:rvachev-up-derivatives}
\end{equation}
from $\up(x)=\Fg(x+1)$ on $[-1,1]$; at a dyadic $x$ of depth $n$ it
vanishes for $r\ge n+1$, and outside the support all derivatives vanish.
The right-hand sides are evaluated by \cref{rf:eq:rvachev-global-fold} and
any finite formula for $\Fb$ at denominator depth at most $n-r$.  Thus
every entry of every dyadic Taylor jet is explicitly rational, and
substituting into \cref{rf:eq:scale-polynomial} or
\cref{rf:eq:rvachev-density-law} makes every coefficient of the scale
polynomials explicit --- for example
$\up''(\tfrac14)=2^{\Tr{2}}\Fg(5)=8\,\tm{2}=-8$ by
\cref{rf:eq:setup-integer-values}.  \Cref{rf:lem:splines-derivative-master} gives the
same derivatives as one Thue--Morse sum without the fold.

\subsection{Iterated primitives}
\label{rf:sec:rvachev-primitives}

For an integer $r\ge1$, repeated integration has an equally simple exact
form:
\begin{align}
 \frac{1}{(r-1)!}\int_{0}^{x}(x-t)^{r-1}\Fg(t)\,\dd t
 &=2^{\binom{r}{2}}\,\Fg\!\left(\frac{x}{2^{r}}\right) &&(x\ge0),
 \label{rf:eq:rvachev-primitive-global}\\
 \frac{1}{(r-1)!}\int_{0}^{x}(x-t)^{r-1}\Fb(t)\,\dd t
 &=2^{\binom{r}{2}}\,\Fb\!\left(\frac{x}{2^{r}}\right) &&(0\le x\le1),
 \label{rf:eq:rvachev-primitive-F}\\
 \frac{1}{(r-1)!}\int_{-1}^{x}(x-t)^{r-1}\up(t)\,\dd t
 &=2^{\binom{r}{2}}\,\Fb\!\left(\frac{x+1}{2^{r}}\right) &&(-1\le x\le1).
 \label{rf:eq:rvachev-primitive-up}
\end{align}

\begin{proof}
The left side of \cref{rf:eq:rvachev-primitive-global} is the $r$-fold
iterated primitive of $\Fg$ with vanishing initial data at $0$ (Cauchy's
formula for repeated integration).  Differentiating the right side $r$ times
with \cref{rf:eq:setup-global-derivative} gives
$2^{\binom{r}{2}}\,2^{\Tr{r}}\,2^{-r^{2}}\Fg(x)=\Fg(x)$, since
$\binom{r}{2}+\Tr{r}-r^{2}=0$; its first $r$ derivatives at $0$ vanish
because $\Fg$ is flat there.  Two functions with the same $r$th derivative
and the same first $r$ initial data coincide.  The second line restricts the
first to $[0,1]$, where $\Fg=\Fb$ and $x/2^{r}\in[0,1]$.  The third uses
$\up(t)=\Fg(t+1)$ and the substitution $t\mapsto t-1$, which moves the lower
endpoint to $0$ and the argument to $(x+1)/2^{r}\in[0,1]$.
\end{proof}

These integrals at dyadic endpoints are therefore rational, with the same
recurrence-free evaluations at a larger denominator depth.

% =====================================================================
\section{The reciprocal-integer-base extension}
\label{rf:sec:base}

The squared-scale mechanism is not confined to base two.  The extension
concerns a \emph{different family of distributions}; it is not a statement
about the Fabius function at non-dyadic rationals, and it does not assert
that the base-$b$ laws share all local analytic properties of $\Fb$.

\subsection{The normalized base-$b$ law and its finite arithmetic}
\label{rf:sec:base-law}

For an integer $b\ge2$ let
\begin{equation}
 X_{b}\DefinitionEquals(b-1)\sum_{j=1}^{\infty}b^{-j}U_{j},\qquad
 \Fb^{(b)}(x)\DefinitionEquals\Prob(X_{b}\le x),
 \label{rf:eq:base-X}
\end{equation}
normalized so that the support is $[0,1]$; $X_{2}=X$.  The same
absolutely convergent series and characteristic-function arguments as in
\cref{rf:lem:setup-normalization} give reflection symmetry
$X_{b}\EqualInLaw1-X_{b}$ and a smooth density, flat at both endpoints.
Write $Z_{b}=2X_{b}-1$, with centred moment generating function
$M_{b}(t)=\prod_{j\ge1}H((b-1)t/b^{j})$, even, and with
\begin{equation}
 \log M_{b}(t)=\sum_{r\ge1}\lambda_{r,b}\,t^{2r},\qquad
 \lambda_{r,b}=\frac{(b-1)^{2r}\,2^{2r-1}B_{2r}}{r\,(2r)!\,(b^{2r}-1)},
 \label{rf:eq:base-lambda}
\end{equation}
by the geometric summation of $\log H$ over the scales $b^{-j}$ exactly as
in \cref{rf:prop:moments-logM}.  The raw coefficients have the positive
finite formula
\begin{equation}
 \boxed{\;
 A^{(b)}_{m}\DefinitionEquals\frac{\E X_{b}^{m}}{m!}
 =(b-1)^{m}\sum_{\mathbf r\in\Comp{m}}\prod_{i=1}^{\ell(\mathbf r)}
 \frac{1}{(r_{i}+1)!\,(b^{R_{i}}-1)}\,,\;}
 \label{rf:eq:base-A}
\end{equation}
by the gap summation \cref{rf:eq:moments-gaps} with $b$ in place of $2$ and
the common scale $(b-1)^{m}$.  Other integer bases do not fill one
consecutive integer block with their subset sums, so a Boolean cube sum
replaces the single Thue--Morse prefix.  For $n\ge1$, $0\le m\le b^{n}$ and
$\delta\in\SetOf{0,1}^{n}$ put
\begin{equation}
 \Delta_{\delta}\DefinitionEquals m-(b-1)\sum_{j=0}^{n-1}\delta_{j}b^{j}.
 \label{rf:eq:base-delta}
\end{equation}

\begin{theorem}[Finite base-$b$ arithmetic]
\label{rf:thm:base-arithmetic}
For these parameters,
\begin{equation}
 \boxed{\;
 \Fb^{(b)}\!\left(\frac{m}{b^{n}}\right)
 =\frac{b^{-\binom{n}{2}}}{(b-1)^{n}}
 \sum_{\substack{\delta\in\SetOf{0,1}^{n}\\ \Delta_{\delta}\ge1}}(-1)^{\AbsoluteValue{\delta}}
 \sum_{l=0}^{n}A^{(b)}_{l}\,\frac{(\Delta_{\delta}-1)^{n-l}}{(n-l)!}\,.\;}
 \label{rf:eq:base-value}
\end{equation}
Consequently $\Fb^{(b)}$ is rational at every base-$b$ rational argument.
\end{theorem}

\begin{proof}
Split $X_{b}$ into its first $n$ coordinates and the independent tail
$b^{-n}X_{b}'$, and apply \cref{rf:lem:master-box} to the first
coordinates, whose weights $(b-1)b^{-j}$, $1\le j\le n$, have product
$(b-1)^{n}b^{-\Tr{n}}$.  Conditional on the tail, the truncated power at
$x=m/b^{n}$ is $b^{-n^{2}}\PositivePart{\Delta_{\delta}-X_{b}'}^{n}$
(indexing the subset by the digits $\delta_{j}$ of the weight
$(b-1)b^{-(n-j)}$).  Since $\Delta_{\delta}$ is an integer and
$X_{b}'\in[0,1]$, this is zero for $\Delta_{\delta}\le0$ and an ordinary
polynomial for $\Delta_{\delta}\ge1$.  Reflection replaces $1-X_{b}'$ by
$X_{b}'$, and expanding $(\Delta_{\delta}-1+X_{b}')^{n}/n!$ gives the inner
sum; the scalar factors combine to $b^{\Tr{n}-n^{2}}(b-1)^{-n}=b^{-\binom{n}{2}}(b-1)^{-n}$.
\end{proof}

\subsection{Parity-reduced extraction at the base $b^{-2}$}
\label{rf:sec:base-extraction}

For $N\ge1$ let $X_{b,N}=(b-1)\sum_{j=1}^{N}b^{-j}U_{j}$ and
$G_{b,N}(x)\DefinitionEquals\Prob(X_{b,N}+\tfrac12b^{-N}\le x)$, the centred
prefix distribution function.  For $N\ge n$ and $A=mb^{N-n}$ the box
formula gives the finite rational expression
\begin{equation}
 G_{b,N}\!\left(\frac{m}{b^{n}}\right)
 =\frac{1}{2^{N}(b-1)^{N}b^{\binom{N}{2}}N!}
 \sum_{\delta\in\SetOf{0,1}^{N}}(-1)^{\AbsoluteValue{\delta}}
 \PositivePart{2A-1-2(b-1)\sum_{j=0}^{N-1}\delta_{j}b^{j}}^{N},
 \label{rf:eq:base-spline}
\end{equation}
in which the positive part is a comparison of explicitly given integers;
at a general rational $x$ the same formula holds with
$2b^{N}x$ in place of $2A$.

\begin{theorem}[Moment-free integer-base extraction]
\label{rf:thm:base-extraction}
For $n\ge1$, $0\le m\le b^{n}$, $d=\Floor{n/2}$, and $R\ge0$,
\begin{equation}
 \boxed{\;
 \Fb^{(b)}\!\left(\frac{m}{b^{n}}\right)
 =\sum_{j=0}^{d}w_{d,j}\!\left(b^{-2}\right)\,G_{b,n+R+j}\!\left(\frac{m}{b^{n}}\right),\;}
 \label{rf:eq:base-extraction}
\end{equation}
with the weights of \cref{rf:eq:extraction-weights-q} at $q=b^{-2}$.
\end{theorem}

\begin{proof}
The self-similarity $X_{b}\EqualInLaw((b-1)U+X_{b}')/b$ gives
\begin{equation}
 \bigl(\Fb^{(b)}\bigr)'(x)=\frac{b}{b-1}\Bigl\{\Fb^{(b)}(bx)-\Fb^{(b)}\bigl(bx-(b-1)\bigr)\Bigr\}.
 \label{rf:eq:base-dilation}
\end{equation}
Iterating this $n$ times writes the $n$th derivative at $x$ as a finite
linear combination of values $\Fb^{(b)}(b^{n}x-k)$ with integral $k$; one
further differentiation and evaluation at $x=m/b^{n}$ involves only
positive-order derivatives of $\Fb^{(b)}$ at integers, which vanish because
integers are endpoints of, or lie outside, the support of the density.
Hence the jet of $\Fb^{(b)}$ at $m/b^{n}$ terminates at order $n$.

Now write $\eta=b^{-N}$ and use the head--tail identity
$\Fb^{(b)}(t)=\E G_{b,N}(t-\tfrac{\eta}{2}Z_{b}')$, with $Z_{b}'$ symmetric
and supported in $[-1,1]$.  Every knot of the centred finite distribution
function is an odd multiple of $\eta/2$, while $m/b^{n}$ is an integer
multiple of $\eta$ for $N\ge n$; the averaging interval of half-width
$\eta/2$ therefore lies in one polynomial cell, and the fact that not every
odd multiple is a knot causes no difficulty.
\Cref{rf:lem:splines-local-inverse} with $W=Z_{b}'$ and the terminating jet
shows that $G_{b,N}(m/b^{n})$ is an exact polynomial of degree at most $d$
in $(\eta/2)^{2}$, hence in $b^{-2N}$, with constant term
$\Fb^{(b)}(m/b^{n})$.  Geometric constant extraction with $q=b^{-2}$ at
the depths $n+R+j$ proves \cref{rf:eq:base-extraction}, and substituting
\cref{rf:eq:base-spline} shows rationality once more.
\end{proof}

For example, the explicit sums give
\[
 \Fb^{(3)}\!\left(\tfrac19\right)=\frac{7}{576},\quad
 \Fb^{(3)}\!\left(\tfrac1{27}\right)=\frac{1}{6912},\quad
 \Fb^{(4)}\!\left(\tfrac1{16}\right)=\frac{1}{240},\quad
 \Fb^{(4)}\!\left(\tfrac1{64}\right)=\frac{1}{51840}.
\]
At $b=2$ the cube sums collapse to the dyadic prefix and every formula of
the preceding sections is recovered.  For $b>2$ some of the highest jets can
vanish for additional geometric reasons, so only the degree \emph{upper
bound} $d$ is asserted, not the exact-degree statement of
\cref{rf:prop:exact-degree}.  The existence of exact finite extraction is
caused by two structures together --- grid-aligned finite-convolution knots
and an even centred tail --- and neither alone would give the degree bound.

% =====================================================================
\section{Worked exact values and the retained verifier}
\label{rf:sec:examples}

\subsection{Universal coefficients}
\label{rf:sec:examples-coefficients}

The first normalized and raw moments, all obtained independently from the
composition, partition, and finite-prefix formulas, are
\begin{center}
\renewcommand{\arraystretch}{1.3}
\small
\begin{tabular}{c@{\quad}r@{\quad}r@{\quad}r@{\quad}r}
\toprule
$k$ & $a_{k}$ & $c_{k}=\E Z^{2k}$ & $e_{k}$ & $\lambda_{k}$\\
\midrule
0 & $1$ & $1$ & $1$ & ---\\
1 & $1/18$ & $1/9$ & $-1/18$ & $1/18$\\
2 & $19/16200$ & $19/675$ & $31/16200$ & $-1/2700$\\
3 & $583/42865200$ & $583/59535$ & $-2347/42865200$ & $1/178605$\\
4 & $132809/1311675120000$ & $132809/32531625$ & $1904369/1311675120000$ & $-1/9639000$\\
\bottomrule
\end{tabular}
\end{center}
These numbers are universal: the dyadic numerator changes only the outer
polynomial evaluation.  The raw coefficients begin
$A_{1}=\tfrac12$, $A_{2}=\tfrac5{36}$, $A_{3}=\tfrac1{36}$,
$A_{4}=\tfrac{143}{32400}$, $A_{5}=\tfrac{19}{32400}$, and the reciprocal
powers $\Fb(2^{-n})=2^{-\binom{n}{2}}A_{n}$ for $n=0,\ldots,7$ are
\[
 1,\ \tfrac12,\ \tfrac5{72},\ \tfrac1{288},\ \tfrac{143}{2073600},\
 \tfrac{19}{33177600},\ \tfrac{1153}{561842749440},\ \tfrac{583}{179789679820800}.
\]

\subsection{Two and three samples}
\label{rf:sec:examples-extraction}

At $x=\tfrac14$ ($n=2$, $d=1$) the explicit spline sums are
$G_{2}(\tfrac14)=\tfrac1{16}$ and $G_{3}(\tfrac14)=\tfrac{13}{192}$ --- the
numerator of the latter is $3^{3}-1^{3}=26$ over $2^{\Tr{3}}3!=384$ --- so
\begin{equation}
 \Fb\!\left(\tfrac14\right)=\frac{-G_{2}(\tfrac14)+4G_{3}(\tfrac14)}{3}=\frac{5}{72},
 \qquad G_{N}\!\left(\tfrac14\right)=\frac5{72}-\frac19\,4^{-N}\ \ (N\ge2).
 \label{rf:eq:examples-quarter}
\end{equation}
Likewise $G_{3}(\tfrac38)=\tfrac{97}{384}$, $G_{4}(\tfrac38)=\tfrac{389}{1536}$,
and $\Fb(\tfrac38)=(4G_{4}-G_{3})/3=\tfrac{73}{288}$; and
$G_{3}(\tfrac18)=\tfrac1{384}$, $G_{4}(\tfrac18)=\tfrac5{1536}$,
$\Fb(\tfrac18)=\tfrac1{288}$ with $G_{N}(\tfrac18)=\tfrac1{288}-\tfrac1{18}4^{-N}$.
At depth four three levels suffice:
\[
 G_{4}\!\left(\tfrac1{16}\right)=\frac{1}{24576},\qquad
 G_{5}\!\left(\tfrac1{16}\right)=\frac{121}{1966080},\qquad
 G_{6}\!\left(\tfrac1{16}\right)=\frac{6331}{94371840},
\]
\begin{equation}
 \Fb\!\left(\tfrac1{16}\right)=\frac{G_{4}-20G_{5}+64G_{6}}{45}=\frac{143}{2073600},
 \label{rf:eq:examples-sixteenth}
\end{equation}
and here the polynomial in $t=4^{-N}$ is genuinely quadratic:
\begin{equation}
 G_{N}\!\left(\tfrac1{16}\right)=\frac{143}{2073600}-\frac{5}{648}\,t+\frac{248}{2025}\,t^{2},\qquad
 G_{N}\!\left(\tfrac3{16}\right)=\frac{46657}{2073600}-\frac{67}{648}\,t-\frac{248}{2025}\,t^{2}
 \quad(N\ge4).
 \label{rf:eq:examples-sixteenth-poly}
\end{equation}
At $x=\tfrac5{16}$, with $G_{4}=\tfrac{1205}{8192}$, $G_{5}=\tfrac{289799}{1966080}$,
$G_{6}=\tfrac{13917509}{94371840}$, the weighted sum
$\tfrac1{45}G_{4}-\tfrac49G_{5}+\tfrac{64}{45}G_{6}$ is exactly
$\tfrac{305857}{2073600}$.  No previously known value of $\Fb$ enters any of
these extractions.

\subsection{A table of exact values}
\label{rf:sec:examples-table}

All entries of \cref{rf:tab:examples-values} were matched by the master,
partition, both binary, determinant, finite-cube and quarter-base formulas
in exact rational arithmetic.  Reflection supplies the remaining values on
each grid: $\Fb((16-m)/16)=1-\Fb(m/16)$, and the bump values follow from
$\up(m/16)=\Fb((16-\AbsoluteValue{m})/16)$, for example
$\up(\tfrac{11}{16})=\Fb(\tfrac5{16})$.

\begin{table}[htbp]
\centering
\caption{Selected exact values, fractions reduced.}
\label{rf:tab:examples-values}
\renewcommand{\arraystretch}{1.3}
\begin{tabular}{@{}rl@{\qquad}rl@{}}
\toprule
$x$ & $\Fb(x)$ & $x$ & $\Fb(x)$\\
\midrule
$1/2$ & $1/2$ & $1/4$ & $5/72$\\
$1/8$ & $1/288$ & $3/8$ & $73/288$\\
$1/16$ & $143/2073600$ & $3/16$ & $46657/2073600$\\
$5/16$ & $305857/2073600$ & $7/16$ & $777743/2073600$\\
$1/32$ & $19/33177600$ & $13/32$ & $10393219/33177600$\\
$21/64$ & $482385079229/2809213747200$ & $37/128$ & $4121049919811/35957935964160$\\
\bottomrule
\end{tabular}
\end{table}
Two further values: $\Fb(1/1024)=134926369/1246394851358539387238350848000
\approx1.082533\cdot10^{-22}$, where the positive formula avoids subtracting
large alternating terms, and $\Fb(173/4096)$ of \cref{rf:eq:binary-large}.
The decimal displays are for orientation; the fractions are the actual
outputs.

\subsection{Operation counts and exact arithmetic}
\label{rf:sec:examples-cost}

For the consecutive-depth quarter formula with $R=0$ the raw integer sums
contain $m(2^{d+1}-1)$ integer-power summands before any use of symmetry or
cancellation identities, and the largest depth is $n+\Floor{n/2}$; reducing
the fraction first, and replacing $x$ by $1-x$ when that is closer to zero,
decreases the count.  The block formula has $\BinaryDigitSum(m)$ outer
terms and needs coefficient indices at most $\Floor{n/2}$ regardless of the
numerator; the telescope has the same number of outer terms with raw
coefficients of index up to $n$.  The signs in $\Spow{N}{A}$ can produce
severe cancellation between very large integers: compute the numerator
exactly, divide only after summation, and combine the interpolation weights
as exact fractions.  The bounded norm of the weight row
(\cref{rf:prop:extraction-norm}) does not make a low-precision evaluation of
the individual sums safe.  None of this asserts that a nonrecursive formula
must outperform a well-designed recurrence.

\subsection{What the retained verifier checks}
\label{rf:sec:examples-verifier}

The retained program \texttt{verify\_recurrence\_free\_dyadic\_values.py}
uses only the Python standard library and \texttt{fractions.Fraction}; every
comparison is an exact equality, with no floating-point tolerance.  It
implements, as separate finite evaluators, the master formula with
composition and with partition coefficients, the raw-moment master formula,
the block formula (bounded and signed), the telescope, the quarter-base row
at consecutive and at arbitrary depths, the finite-cube formula, the
bordered determinant (by exact Bareiss elimination), the certified rounding,
the half-base row with arbitrary shift $\alpha$ and its closed form, the
direct density extraction, and the integer-base formulas; the only
recursive object, the triangular moment recurrence, is isolated and used
solely as a cross-check.  Its recorded run (\texttt{verification\_results.json})
performs the checks of \cref{rf:tab:examples-checks}.

\begin{table}[htbp]
\centering
\caption{Exact checks performed by the retained verifier at maximal depth 8.}
\label{rf:tab:examples-checks}
\small
\begin{tabularx}{\textwidth}{@{}Xr@{}}
\toprule
Check & Cases\\
\midrule
Coefficients $a_{k}$ by compositions, partitions and the isolated oracle, $0\le k\le12$; recovery from finite prefixes at two starting depths, multinomial against power-sum prefix coefficients, both reciprocal expansions, $(-1)^{k}e_{k}>0$ and $\sum_{j}a_{j}e_{k-j}=[k=0]$, $0\le k\le8$; $\lambda_{r}$ by the Bernoulli and the Bernoulli-free formula, $r\le8$; raw $A_{m}$ by compositions, partitions and the dictionary \cref{rf:eq:setup-raw-centred}, $m\le8$; every printed constant & 39\\
Weight rows $d\le8$: product, $q$-Pochhammer, Gaussian-binomial and integer forms agree; annihilation identities; absolute row sum and generating polynomial at three bases; the bound $<2$ & 9 rows\\
Every representation $m/2^{n}$, $0\le n\le8$: master (three coefficient routes), block, signed block, telescope, quarter-base, determinant (and finite-cube for $n\le6$) agree; reflection; refinement $(m,n)\mapsto(2m,n+1)$; integrality against $D_{n}$ and $D'_{n}$ & 1{,}033\\
Shifted-depth and arbitrary-depth extraction at selected points & 84\\
Signed extension $\Fg(a/2^{n})$ by blocks against the fold, $0\le a\le3\cdot2^{n}$, $n\le6$; integer and half-integer values & 388\\
Exact shrinking-cell identity at five offsets, grids through depth 5, three depths; scale polynomials through depth 6 at three depths; exact degree at every reduced interior point; the four printed polynomials & 1{,}005\\
Dyadic derivatives by the differentiated master formula against the dilation law, orders $0\le r\le n+1$, $n\le5$ & 434\\
Certified rounding, depths 1--3 & 17\\
Half-base row with $\alpha\in\SetOf{0,\tfrac12,1,-2,\tfrac23}$, and its closed form & 65\\
Rvachev values by direct density extraction (two starting depths), the finite fold identity, the density law \cref{rf:eq:rvachev-density-law}, $\up''(\tfrac14)=-8$ & 126\\
Odd-depth and reciprocal-power formulas; the eight anchor values & 22\\
Integer bases $b=2,3,4,5$: extraction against the finite arithmetic formula and at a shifted depth, reflection, cube-sum evaluations & 235\\
\bottomrule
\end{tabularx}
\end{table}

The grid cases are numerator--depth pairs, not distinct rational points:
unreduced representations are intentionally included.  To reproduce the
recorded run:
\begin{lstlisting}[language=bash]
python verify_recurrence_free_dyadic_values.py --max-depth 8 --json verification_results.json
\end{lstlisting}
The exhaustive sums grow exponentially with the requested depth.  The tests
check the implementations and the indexing; they are not a substitute for
the proofs, nor a claim that any of this has been formalized in Lean.

% =====================================================================
\section{How the alternatives compare, and what remains}
\label{rf:sec:conclusion}

\subsection{Closed form versus evaluation algorithm}
\label{rf:sec:conclusion-compare}

An explicit finite formula may be evaluated by loops, memoization, or an
exact linear-algebra algorithm without ceasing to be a closed form; what is
excluded is a definition that leaves an auxiliary constant specified only
through earlier unknown constants.  Storing previously evaluated $a_{k}$ is
optional caching, whereas defining $a_{k}$ only by
\cref{rf:eq:moments-recurrence} would fail the criterion.

The families are algebraically equivalent descriptions but not identical
algorithms.  The composition sum has $2^{k-1}$ positive terms, the partition
sum one signed term per integer partition, the finite-prefix representation
only weak compositions and rational Gaussian weights.  Once coefficients are
available, the block formula has at most $(d+1)(n-d+1)$ polynomial terms.
The quarter-base double sum is the most direct answer with no named
coefficient sequence, at the price of $m(2^{d+1}-1)$ signed powers and
enormous intermediate integers.  The determinant supplies a common
denominator immediately, but its permutation expansion is not the preferred
numerical algorithm.  Large powers and cancellation, bit lengths of rational
intermediates, coefficient reuse, and whether one value or a whole grid is
wanted all matter more than the number of summands.

\subsection{Further deductions that need no new conjecture}
\label{rf:sec:conclusion-further}

A finite subset of the cutoff coefficients can vanish at special points;
whenever the nonzero orders are known one may interpolate only those
monomials rather than the complete degree space.  The exact-degree result
guarantees that the top order does not vanish at reduced points but does not
exclude gaps at lower orders, which suggests point-adapted extraction rows
whose validity is proved by the corresponding finite moment equations.  The
local inverse lemma applies more generally whenever a smooth law is an
independent sum of a finite spline head and a scaled tail, the tail's
support fits in one polynomial cell at the evaluation point, and the
relevant jet of the limit terminates: symmetry removes the odd cutoff
powers, other coefficient zeros can remove further powers, and the search
for additional exact formulas separates into a geometric question about
knot cells and an algebraic one about the nonzero jet orders.  The
integer-base theorem is one realization of that separation.  Open beyond
this document: a formalization of the general polynomial law
(\cref{rf:sec:register}), and the sharp denominator arithmetic of
\cite{arias2018} against the elementary bounds $D_{n}$, $D'_{n}$.

\subsection{Conclusion}
\label{rf:sec:conclusion-conclusion}

Every dyadic value in question admits several closed forms in the strict
finite sense.  The hidden moment recurrences of the classical formula can
be eliminated by explicit multiplicity partitions, by positive ordered
compositions, by finite-prefix interpolation, or by an integer determinant.
Binary digits provide two finite compressions with no long signed prefix.
Most importantly, centred finite convolutions have an exact polynomial
dependence on the squared tail scale at dyadic arguments; removing that
polynomial by geometric interpolation gives \cref{rf:eq:main-quarter}, a
moment-free $q=\tfrac14$ formula requiring only $\Floor{n/2}+1$ spline
levels, and the shrinking-cell theorem explains why the quarter base
replaces the half base and why that number of samples is the least
possible for a grid-wide rule.  Reflection and folding cover every Rvachev
and signed Fabius value, and the same finite algebra extends to all
derivatives, iterated primitives, and reciprocal-integer-base laws.
Infinite products and probability models appear in the proofs; neither
infinite limits nor auxiliary recurrences survive in the evaluation
formulas.

\clearpage
\appendix

% =====================================================================
\section{Compact exact evaluators}
\label{rf:app:code}

The following self-contained implementations use only integer powers,
factorials, binary digit counts, and exact fractions.  The first is the
principal formula \cref{rf:eq:main-quarter}; the second is the block formula
\cref{rf:eq:binary-block} with the composition coefficients
\cref{rf:eq:moments-a-composition}, valid for the signed extension at every
nonnegative numerator.  Python 3.10 or later suffices; pass exact rational
inputs such as \texttt{Fraction(173, 4096)} rather than decimal floats,
whose exact binary value may have a much larger denominator than intended.

\begin{lstlisting}[language=Python]
from fractions import Fraction
from functools import lru_cache
from math import factorial, prod


def fabius_quarter(m: int, n: int) -> Fraction:
    """Exact bounded F(m / 2**n) for 0 <= m <= 2**n by the quarter-base row."""
    if n < 0 or not 0 <= m <= 2**n:
        raise ValueError("require n >= 0 and 0 <= m <= 2**n")
    while n > 0 and m % 2 == 0:      # reduce: fewer levels
        m //= 2; n -= 1
    if n == 0:
        return Fraction(m)
    d = n // 2
    total = Fraction(0)
    for j in range(d + 1):
        weight = Fraction((-1)**(d - j) * 4**(j*(j+1)//2),
                          prod(4**r - 1 for r in range(1, j+1))
                          * prod(4**r - 1 for r in range(1, d-j+1)))
        N, A = n + j, m * 2**j
        power_sum = sum((-1)**h.bit_count() * (2*A - 2*h - 1)**N
                        for h in range(A))
        total += weight * Fraction(power_sum, 2**(N*(N+1)//2) * factorial(N))
    return total


def compositions(k: int):
    if k == 0:
        yield (); return
    for mask in range(1 << (k - 1)):
        cuts = [0] + [j for j in range(1, k) if (mask >> (j-1)) & 1] + [k]
        yield tuple(v - u for u, v in zip(cuts, cuts[1:]))


@lru_cache(None)
def coefficient(k: int) -> Fraction:
    """a_k = [t^(2k)] M(t) by the positive composition formula."""
    total = Fraction(0)
    for parts in compositions(k):
        term, suffix = Fraction(1), k
        for r in parts:
            term /= factorial(2*r + 1) * (4**suffix - 1); suffix -= r
        total += term
    return total


def fabius_global_blocks(a: int, n: int) -> Fraction:
    """Signed global F_gl(a / 2**n), every a >= 0, by the binary block formula."""
    if a < 0 or n < 0:
        raise ValueError("a and n must be nonnegative")
    total = Fraction(0)
    for i in range(n + 1):
        if not (a >> i) & 1:
            continue
        sign = -1 if (a >> (i + 1)).bit_count() % 2 else 1
        K = (1 << i) + 2 * (a % (1 << i))
        inner = sum((coefficient(k) * 4**(i*k)
                     * Fraction(K**(n-i-2*k), factorial(n-i-2*k))
                     for k in range((n-i)//2 + 1)), Fraction(0))
        total += sign * 2**(i*(i+1)//2) * inner
    return total / 2**(n*(n+1)//2)


def rvachev_up(a: int, n: int) -> Fraction:
    """Exact up(a / 2**n) on the whole dyadic grid."""
    return Fraction(0) if abs(a) > 2**n else fabius_quarter(2**n - abs(a), n)


assert fabius_quarter(1, 4) == Fraction(143, 2073600)
assert fabius_global_blocks(173, 12) == fabius_quarter(173, 12)
assert fabius_global_blocks(3, 0) == -1 and rvachev_up(3, 2) == Fraction(5, 72)
\end{lstlisting}

% =====================================================================
\section{A direct finite definition of every index set}
\label{rf:app:indices}

To remove any ambiguity about the word ``finite'', all auxiliary indexing
reduces to bounded integer tuples.  A composition of $k\ge1$ is a length
$1\le\ell\le k$ and a tuple $(r_{1},\ldots,r_{\ell})\in\SetOf{1,\ldots,k}^{\ell}$
with $r_{1}+\cdots+r_{\ell}=k$; the empty tuple is the only composition of
$0$.  A multiplicity vector of weight $k$ lies in
$\prod_{r=1}^{k}\SetOf{0,1,\ldots,\Floor{k/r}}$ and satisfies
$\sum_{r}rv_{r}=k$; a raw profile of weight $m$ adds
$v_{0}\in\SetOf{0,\ldots,m}$ with $v_{0}+2\sum_{r}rv_{r}=m$.  Binomial
coefficients are factorial quotients in their stated ranges; Bernoulli
numbers are the finite sums \cref{rf:eq:moments-bernoulli-finite}; every
Thue--Morse sign is the digit product \cref{rf:eq:intro-tm-digits}; every
determinant is the permutation sum \cref{rf:eq:determinant-leibniz}; every
$q$-Pochhammer symbol in a value formula is the finite product
\cref{rf:eq:intro-q}; the digit data \cref{rf:eq:binary-parameters} and
\cref{rf:eq:binary-suffix} are floor and remainder operations.  There is no
unstated recursive sequence anywhere in the evaluation layer.

% =====================================================================
\section{Notation}
\label{rf:sec:notation}

Every symbol is normative from the corpus notation catalogue where an entry
exists (\texttt{rvachev-up}, \texttt{fabius-bounded},
\texttt{fabius-global}, \texttt{thue-morse}, \texttt{q-pochhammer},
\texttt{gaussian-binomial}); the remaining symbols are document-local and
listed here.

\begin{center}
\renewcommand{\arraystretch}{1.2}
\small
\begin{tabularx}{\textwidth}{@{}lX@{}}
\toprule
Symbol & Meaning\\
\midrule
$\Fb$, $\up$, $\Fg$ & bounded distribution function of $X$; density of $Z=2X-1$; signed global extension \cref{rf:eq:setup-global}\\
$X$, $Z$; $U_{j}$, $V_{j}$ & the random sums \cref{rf:eq:setup-X,rf:eq:setup-Z}; uniforms on $[0,1]$ and $[-1,1]$\\
$X_{N}$, $Z_{N}$, $Y_{N}$ & uncentred and centred prefixes; $Y_{N}=X_{N}+2^{-N-1}$\\
$H_{N}$, $G_{N}$, $G_{N,\alpha}$, $u_{N}$ & distribution functions of $X_{N}$, $Y_{N}$, $X_{N}+\alpha2^{-N}$; density of $Z_{N}$\\
$H$, $M$, $M_{N}$, $M_{X}$ & $\sinh t/t$; moment generating functions of $Z$, $Z_{N}$, $X$\\
$a_{k}$, $c_{k}$, $A_{m}$ & $\coeff{t}{2k}M$, $\E Z^{2k}$, $\E X^{m}/m!$\\
$\lambda_{r}$, $\beta^{\ast}_{r}$, $e_{k}$ & coefficients of $\log M$, of $\log M_{X}-z/2$, of $1/M$\\
$\Phi_{r}(v)$, $C_{n,N}(A)$, $a_{k,N}$ & $\coeff{t}{r}\EulerE^{vt}M$; $\E(A+Z_{N})^{n}/n!$; $\coeff{t}{2k}M_{N}$\\
$\Tr{r}$, $d$, $\tm{h}$ & $r(r+1)/2$; $\Floor{n/2}$; Thue--Morse sign $(-1)^{\BinaryDigitSum(h)}$\\
$\Spow{N}{A}$, $t_{j}$, $L_{d}$, $B_{n,m}$, $D_{n}$, $D'_{n}$ & integer power sum \cref{rf:eq:intro-power-sum}; determinant data \cref{rf:sec:determinant}\\
$\Comp{k}$, $R_{i}$, $\Prof{k}$ & compositions, suffix sums, multiplicity profiles \cref{rf:sec:intro-conventions}\\
$w_{d,j}(q)$, $w_{d,j}$ & geometric extraction weights \cref{rf:eq:extraction-weights-q}; at $q=\tfrac14$\\
$P_{x}$, $Q_{x}$ & dyadic Taylor polynomial; scale polynomial \cref{rf:eq:scale-polynomial}\\
$J(m)$, $p_{j}$, $r_{j}$, $K_{j}$, $\eta_{j}$; $b_{i}$, $s_{i}$ & block data \cref{rf:eq:binary-parameters}; telescope data \cref{rf:eq:binary-suffix}\\
$X_{b}$, $\Fb^{(b)}$, $G_{b,N}$, $A^{(b)}_{m}$, $\lambda_{r,b}$ & the base-$b$ objects of \cref{rf:sec:base}\\
$\amvol{r}$ & the canonical volume's reciprocal-product coefficients, $=(-1)^{r}e_{r}$\\
\bottomrule
\end{tabularx}
\end{center}

% =====================================================================
\section{Formalization register}
\label{rf:sec:register}

The results of this document are human proofs.  The repository's Lean
development \cite{proveit-primary} contains counterparts for several of the
ingredients, which the canonical volume's register \cite{proveit-due}
already lists in the density normalization; the rows below state the
correspondence for the objects as they appear here.  Every declaration
name was checked to resolve in the repository on the day of consolidation.
Statuses: \textsc{crosswalked}, a declaration states the same fact (with
the translation stated); \textsc{instance}, kernel-verified for one point
or one mode; \textsc{absent}, no declaration found.

\begin{center}
\small
\begin{longtable}{@{}>{\raggedright\arraybackslash}p{0.30\textwidth}>{\raggedright\arraybackslash}p{0.46\textwidth}>{\raggedright\arraybackslash}p{0.18\textwidth}@{}}
\caption{Formalization register.}\label{rf:tab:register}\\
\toprule
Result here & Lean declaration(s) & Status\\
\midrule
\endfirsthead
\toprule
Result here & Lean declaration(s) & Status\\
\midrule
\endhead
\cref{rf:eq:splines-G-integer}: the centred finite spline as a Thue--Morse truncated-power sum
& \texttt{Fabius.\allowbreak{}fabiusUniformSpline} (\texttt{FabiusUniformSpline}), defined as that sum in the centred-CDF variable; \texttt{Fabius.\allowbreak{}ProbabilityRepresentation.\allowbreak fabiusUniformSpline\_\allowbreak{}eq\_\allowbreak{}centeredPartialCDF}
& \textsc{crosswalked}\\[0.3em]
\cref{rf:eq:setup-Z}, \cref{rf:lem:setup-normalization}: $\up$ is the density of $\sum_{k}2^{-k}V_{k}$
& \texttt{Fabius.\allowbreak{}ProbabilityRepresentation.\allowbreak rvachevUp\_\allowbreak{}eq\_\allowbreak{}weightedSum\_\allowbreak{}probability}, \texttt{\ldots geometricUniformDensity\_\allowbreak{}one\_\allowbreak{}half\_\allowbreak{}eq\_\allowbreak{}rvachevUp}
& \textsc{crosswalked}\\[0.3em]
\cref{rf:eq:moments-e-partition}, \cref{rf:eq:moments-a-partition}: the reciprocal and moment coefficients in Bell form
& \texttt{Fabius.\allowbreak{}rvachevReciprocalMomentRat}, \texttt{\ldots\_\allowbreak{}eq\_\allowbreak{}completeBellPolynomial}, \texttt{Fabius.\allowbreak{}binomialConv\_\allowbreak{}rvachevRawMomentRat\_\allowbreak{}reciprocal} (\texttt{RvachevMomentAppell}): $\amvol{r}=(-1)^{r}e_{r}$ is Lean-computable
& \textsc{crosswalked}; positivity \cref{rf:lem:splines-e-sign}: \textsc{absent}\\[0.3em]
\cref{rf:lem:splines-local-inverse} at scale $1$
& \texttt{Fabius.\allowbreak{}rvachevDeconvolvedPolynomial}, \texttt{Fabius.\allowbreak{}rvachevAppellPolynomialRat}, \texttt{Fabius.\allowbreak{}integral\_\allowbreak{}eval\_\allowbreak{}rvachevAppellPolynomial\_\allowbreak{}sub\_\allowbreak{}mul\_\allowbreak{}rvachev} (\texttt{RvachevMomentAppell})
& \textsc{crosswalked} at scale $1$; the dilation and the cell identity: \textsc{absent}\\[0.3em]
\cref{rf:lem:setup-dilation}: dilation law and the terminating dyadic jet
& \texttt{Fabius.\allowbreak{}deriv\_\allowbreak{}rvachevUp}, \texttt{Fabius.\allowbreak{}iteratedDeriv\_\allowbreak{}rvachevUp\_\allowbreak{}eq\_\allowbreak{}thueMorse\_\allowbreak{}sum} (\texttt{Differential}); \texttt{Fabius.\allowbreak{}iteratedDeriv\_\allowbreak{}rvachevUp\_\allowbreak{}dyadic\_\allowbreak{}eq\_\allowbreak{}zero}, \texttt{Fabius.\allowbreak{}dyadic\_\allowbreak{}depth\_\allowbreak{}eq\_\allowbreak{}max\_\allowbreak{}nonzero\_\allowbreak{}iteratedDeriv} (\texttt{DyadicDerivativeFiltration}); \texttt{Fabius.\allowbreak{}iteratedDeriv\_\allowbreak{}fabiusReal\_\allowbreak{}dyadicGrid\_\allowbreak{}eq\_\allowbreak{}ite}, \texttt{Fabius.\allowbreak{}abs\_\allowbreak{}iteratedDeriv\_\allowbreak{}fabiusReal\_\allowbreak{}dyadicGrid\_\allowbreak{}of\_\allowbreak{}odd} (\texttt{BoundedDerivatives})
& \textsc{crosswalked} (the top derivative $2^{\Tr{n}}$ at odd numerators is \cref{rf:prop:exact-degree}'s nonvanishing)\\[0.3em]
\cref{rf:eq:quarter-cell}: the cell identity at $x=\tfrac14$
& \texttt{Fabius.\allowbreak{}fabiusUniformSpline\_\allowbreak{}quarter\_\allowbreak{}eq\_\allowbreak{}quadratic}, \texttt{Fabius.\allowbreak{}reportFiniteFabiusApproximant\_\allowbreak{}quarter\_\allowbreak{}eq\_\allowbreak{}quadratic} (\texttt{QuarterSplineLocalPolynomial}): $\tfrac5{72}+z+4z^{2}-\tfrac49\,4^{-n}$ in the volume's level index $n=N-1$
& \textsc{instance}; \cref{rf:thm:local-cell,rf:cor:scale-polynomial} at general depth: \textsc{absent}\\[0.3em]
\cref{rf:lem:extraction-weights}: the weights, their closed forms, row polynomial and moment identities
& \texttt{Fabius.\allowbreak{}geometricLagrangeWeight}, \texttt{\ldots\_\allowbreak{}eq\_\allowbreak{}product} (\texttt{GeometricLagrangeWeights}); \texttt{Fabius.\allowbreak{}quarterLagrangeWeight\_\allowbreak{}eq\_\allowbreak{}qPochhammer}, \texttt{\ldots\_\allowbreak{}eq\_\allowbreak{}qBinomial} (\texttt{GeometricLagrangeQBinomial}); \texttt{Fabius.\allowbreak{}sum\_\allowbreak{}quarterLagrangeWeight\_\allowbreak{}eq\_\allowbreak{}one}, \texttt{\ldots\_\allowbreak{}mul\_\allowbreak{}pow\_\allowbreak{}eq\_\allowbreak{}zero}, \texttt{Fabius.\allowbreak{}quarterLagrangeWeightPolynomial\_\allowbreak{}eq\_\allowbreak{}forwardRichardson}
& \textsc{crosswalked}\\[0.3em]
\cref{rf:thm:main-quarter}, \cref{rf:thm:extraction-consecutive}: the row applied to the spline values
& \texttt{Fabius.\allowbreak{}quarterLagrange\_\allowbreak{}rvachevFourierProduct\_\allowbreak{}eq\_\allowbreak{}gaussian\_\allowbreak{}tsum} (\texttt{RvachevQBinomialFilter}) applies the same row to the infinite product on the Fourier side
& neighbouring only; the theorem: \textsc{absent}, pending the law\\[0.3em]
\cref{rf:thm:master}, \cref{rf:thm:moments-composition}, \cref{rf:thm:moments-partition}, \cref{rf:thm:binary-block}, \cref{rf:thm:binary-telescope}, \cref{rf:thm:moment-extraction}, \cref{rf:thm:finite-cube}, \cref{rf:thm:determinant}, \cref{rf:cor:rounding}, \cref{rf:prop:halfbase}, \cref{rf:thm:up-quarter}, \cref{rf:sec:rvachev-primitives}, \cref{rf:sec:base}
& --- & \textsc{absent}\\
\bottomrule
\end{longtable}
\end{center}

A formalization of the general polynomial law would follow the five-step
plan recorded in the canonical volume (finite density and knot geometry,
head--tail factorization at finite level, dilated deconvolution and the
cell identity, assembly with the jet termination); the master identity and
the coefficient formulas of \cref{rf:sec:master,rf:sec:moments} need only
finite algebra over the existing \texttt{RvachevMomentAppell} coefficients
and the box formula \cref{rf:lem:master-box}.

% =====================================================================
\section{Provenance of the consolidation}
\label{rf:sec:provenance}

\subsection{What was merged}

Three independently written articles arrived on 2026-09-04 and were filed
the same day, as separate members beside the canonical volume, in the
subgroup \path{inverse-and-sampling/dyadic-up-extraction/}.  All three
answer the same question --- exact evaluation of $\Fb(m/2^{n})$ and
$\up(m/2^{n})$ with no auxiliary quantity defined by a moment recurrence or
a limiting process --- and all three reach the quarter-base formula through
the same mechanism.  They were merged editorially into this document on
2026-09-04: one statement of each result, every proof completed, the four
letters for the numerator, the three names for the moment coefficients, the
three names for the reciprocal coefficients, and the two spline indexings
reconciled once, and every printed rational checked by the retained
verifier.  The three directories and their retained arrival PDFs were
deleted after a residue audit; git history is the archive.

\subsection{The absorbed sources and what each contributed}

\begin{description}[leftmargin=1.5em,style=nextline]
\item[S1: \emph{Recurrence-Free Closed Forms for Dyadic Fabius and Rvachev
Values}] (\texttt{fabius\_dyadic\_closed\_forms/}, 1{,}441 lines, 22 pp.).
The Taylor-integral proof of the master identity for the signed extension
(\cref{rf:rmk:master-taylor}); the block formula in integer form and its
validity at every numerator (\cref{rf:thm:binary-block}); the Bernoulli-free
cumulant formula (\cref{rf:lem:moments-lambda-nob}); the finite-cube value
formula and the arbitrary-prefix weights (\cref{rf:thm:finite-cube}); the
$(d+2)$-bordered determinant (\cref{rf:rmk:determinant-alt}); the Prouhet
route to the degree drop (\cref{rf:rmk:splines-prouhet}); the density
splines $S_{N}=G_{N-1}$; the examples $\Fb(1/1024)$, $\Fb(173/4096)$ and the
three-sample extraction at $\tfrac5{16}$.
\item[S2: \emph{Recurrence-Free Dyadic Values of the Fabius and Rvachev
Functions}] (\texttt{fabius\_dyadic\_closed\_forms-2/}, 1{,}456 lines,
26 pp.).  The organization of \cref{rf:sec:intro,rf:sec:setup,rf:sec:master};
the exact shrinking-cell theorem (\cref{rf:thm:local-cell}); the exact
degree and the minimal-sample theorem
(\cref{rf:prop:exact-degree,rf:thm:extraction-minimal}); the bound $<2$ on
the row (\cref{rf:prop:extraction-norm}); the finite-prefix coefficients and
their extraction (\cref{rf:sec:prefix}); the differentiated master formula
(\cref{rf:lem:splines-derivative-master}); the composition denominator bound,
the prefix error and the certified rounding
(\cref{rf:prop:determinant-denominator,rf:cor:rounding}); the integer-base
cumulants and the centred base-$b$ prefix; the example $21/64$ and the
quadratic scale polynomials at $\tfrac1{16}$, $\tfrac3{16}$.
\item[S3: \emph{Recurrence-Free Dyadic Formulae for the Fabius and Rvachev
Functions}] (\texttt{fabius\_rvachev\_recurrence\_free\_closed\_forms/},
1{,}681 lines, 23 pp.).  The raw-moment master identity and the raw
composition and partition formulas
(\cref{rf:eq:master-raw,rf:thm:moments-composition,rf:thm:moments-partition});
the odd-depth formula \cref{rf:eq:moments-odd-depth}; the binary telescope
(\cref{rf:thm:binary-telescope}); the $(d+1)$-bordered determinant and
$D_{n}$ (\cref{rf:thm:determinant}); the general finite-jet deconvolution
lemma (\cref{rf:lem:splines-local-inverse}); the row's generating
polynomial and the numerical limits of the row norm; the half-base identity
with arbitrary shift (\cref{rf:prop:halfbase}); the direct density
extraction (\cref{rf:thm:up-quarter}); the iterated primitives
(\cref{rf:sec:rvachev-primitives}); the finite base-$b$ arithmetic and
extraction (\cref{rf:thm:base-arithmetic,rf:thm:base-extraction}); the
index-set appendix.
\end{description}

\subsection{Conventions reconciled}

S1 wrote the numerator as $a$, the normalized moments as $b_{k}$, the
reciprocal coefficients as $\alpha_{r}$, the signed extension as $\Phi$, and
worked with density splines $S_{N}(x)=p_{N}(x-1)$, so that its levels
$n+1,\ldots,n+1+d$ are the levels $n,\ldots,n+d$ of $G_{N}$ here; S2 wrote
$m$, $a_{k}$, $e_{k}$, $\Theta$, $G_{N}$, which this document adopts; S3
wrote $a$, $C_{j}$, $R_{j}$, $\mathcal F$, $\mathsf H_{N}=G_{N}$ and the
cumulant coefficients $\lambda_{r}$ (S1's $\beta_{r}$), adopted here, with
S3's raw $\beta_{r}$ renamed $\beta^{\ast}_{r}$.  S1's finite Bernoulli sum
and S2/S3's are the same formula.  The three common-denominator statements
were two: $D_{n}$ (S1, S3) and $D'_{n}$ (S2), both kept.  Where the sources
stated a result without proof --- the Bernoulli-free cumulant formula, the
coefficient-only prefix formula, the generating polynomial of the row, the
primitives --- a proof was written; where a proof was sketched --- the
integer-base extraction in S2, the raw-to-centred dictionary, the
determinant's Schur step --- it was completed.  No mathematical claim of
any source was found to be false; the only corrections were indexing
conventions, which are recorded above.

\subsection{Status of the claims}

Every theorem, proposition, lemma and corollary carries a proof.  Every
displayed rational number, weight row, polynomial coefficient and example
value is asserted by the retained verifier
(\cref{rf:sec:examples-verifier}), whose recorded run passed with the
counts of \cref{rf:tab:examples-checks}.  The verifier merges the three
retired programs of the sources (\texttt{verify\_formulas.py},
\texttt{verify\_closed\_forms.py}, \texttt{verify.py}) and adds the checks
introduced by the consolidation: the raw-to-centred dictionary, the
Bernoulli-free cumulants, the second reciprocal expansion, the arbitrary-depth
weights, the closed half-base form, the density law, the derivative formula
against the dilation law, and every constant printed here.  No Lean
formalization of the new material is claimed; \cref{rf:sec:register} is
the boundary.

% =====================================================================
\begin{thebibliography}{99}
\setlength{\itemsep}{3pt}

\bibitem{proveit-primary}
V. Reshetnikov and contributors,
\emph{The Fabius Function and Rvachev's Up-Function: Exact arithmetic,
probability, Fourier analysis, and corrected sharp asymptotics},
primary exposition of the \texttt{ProveIt} repository,
\path{Analysis/FabiusFunction/docs/Fabius_Function_and_Rvachev_Up/}, together
with the Lean development under \path{Analysis/FabiusFunction/}.
\url{https://github.com/VladimirReshetnikov/ProveIt}.

\bibitem{proveit-due}
V. Reshetnikov,
\emph{Exact Dyadic Extraction of Rvachev's Up-Function from Finite
Sinc-Product Splines}, canonical consolidated volume of the subgroup
\path{drafts/inverse-and-sampling/dyadic-up-extraction/Dyadic_Up_Extraction/},
3 September 2026.

\bibitem{proveit-exponents}
V. Reshetnikov,
\emph{Geometric $q$-Fabius Frontiers}, consolidated volume
\path{drafts/exponents-and-q-series/geometric_q_fabius_frontiers/}, Part III
(the finite sinc-product splines and the quarter-base Richardson combination).

\bibitem{arias1982}
J. Arias de Reyna,
\emph{An infinitely differentiable function with compact support: Definition
and properties}, English translation (2017) of the paper in
\emph{Rev. Real Acad. Ciencias Exactas F\'{\i}s. Natur. Madrid} \textbf{76}
(1982), 21--38.  \url{https://arxiv.org/abs/1702.05442}.

\bibitem{arias2018}
J. Arias de Reyna,
\emph{Arithmetic of the Fabius function},
\emph{Integers} \textbf{18} (2018), Paper A51, 20 pp.
\url{https://math.colgate.edu/~integers/s51/s51.pdf};
preprint \url{https://arxiv.org/abs/1702.06487}.

\bibitem{haugland}
J. K. Haugland,
\emph{Evaluating the Fabius function},
arXiv:1609.07999, version 2 (2020; first version 2016).
\url{https://arxiv.org/abs/1609.07999}.

\bibitem{reshetnikov2019}
V. Reshetnikov,
\emph{Conjectured formula for the Fabius function},
Mathematics Stack Exchange, question posted 5 July 2019.
\url{https://math.stackexchange.com/questions/3283519/conjectured-formula-for-the-fabius-function}.

\bibitem{pinelis2020}
I. Pinelis,
answer to \emph{A non-recursive, explicit formula for the Fabius function},
MathOverflow, January 2020.
\url{https://mathoverflow.net/questions/351004/a-non-recursive-explicit-formula-for-the-fabius-function}.

\bibitem{dlmf}
NIST Digital Library of Mathematical Functions,
\S4.22 (infinite products and partial fractions, Eq.~4.22.1),
\S17.2 ($q$-calculus), \S24.2 (Bernoulli numbers).
\url{https://dlmf.nist.gov/}.

\end{thebibliography}
\end{document}
